\documentclass[12pt]{article}
\usepackage{geometry}
\geometry{left=1in,right=0.75in,top=1in,bottom=1in}

\usepackage[utf8]{inputenc}
\usepackage[T1]{fontenc}

% ---- math / theorem envs ----
\usepackage{amsmath,amssymb,bm}
\usepackage{amsthm,amsfonts}

% ---- figures / tables ----
\usepackage{graphicx}
\usepackage{booktabs}
\usepackage{array}
\usepackage{tabularx}
\usepackage{subcaption}

\newcolumntype{L}{>{\raggedright\arraybackslash}X}
\newcolumntype{C}{>{\centering\arraybackslash}X}

% ---- algorithms (optional) ----
\usepackage{algorithm}
\usepackage{algpseudocode}

% ---- tikz/pgfplots (optional) ----
\usepackage{tikz}
\usepackage{pgfplots}
\pgfplotsset{compat=1.18}

% ---- hyperlinks (article MUST load hyperref explicitly) ----
\usepackage[colorlinks=true,linkcolor=blue,citecolor=blue,urlcolor=blue]{hyperref}

% ---- author/affiliation support for article ----
\usepackage{authblk}

% ---- theorem-like environments ----
\newtheorem{definition}{Definition}
\newtheorem{remark}{Remark}
\newtheorem{assumption}{Assumption}
\newtheorem{proposition}{Proposition}
\newtheorem{theorem}{Theorem}
\newtheorem{lemma}{Lemma}
\newtheorem{corollary}{Corollary}

% ---- keywords: store in preamble, print after abstract ----
\makeatletter
\newcommand{\keywords}[1]{\gdef\@keywords{#1}}
\newcommand{\printkeywords}{%
	\par\vspace{0.5em}%
	\noindent\textbf{Keywords:} \@keywords\par
}
\makeatother

% ---- title & authors ----
\title{Reliability-Aware Multi-objective Regime Analysis of Lunar Construction Logistics with Space-Elevator Transshipment}

\author[1]{Qiurui Liu}
\author[1]{Jiaqi Huang}
\author[2,*]{Yong Wang}

\affil[1]{School of International Exchange, Guangzhou Maritime University, Guangzhou 510725, China}
\affil[2]{School of Arts and Sciences, Guangzhou Maritime University, Guangzhou 510725, China}
\affil[*]{Corresponding author: \texttt{p117646@siswa.ukm.edu.my}}
\date{}

\begin{document}
	\maketitle
	
	\begin{abstract}
		Large-scale lunar construction will require sustained logistics over long planning horizons, yet the competitiveness of alternative transport architectures remains unclear once throughput limits, operational interruptions, and transshipment constraints are considered jointly. In particular, although space-elevator-assisted transport is often discussed as a high-capacity alternative to repeated direct Earth--Moon launches, the conditions under which it becomes systemically competitive remain insufficiently understood.
		
		Here we develop a reliability-aware, campaign-scale framework for comparing three classes of lunar logistics architectures: direct rocket delivery, two-stage space-elevator-assisted transshipment, and hybrid strategies that combine both pathways. Space-elevator reliability is represented by a continuous-time Markov chain and incorporated into campaign-scale planning through an availability-adjusted effective-throughput factor. Architecture evaluation is formulated as a multi-objective optimization problem involving construction completion time, a normalized cost index, and an Earth-launch-intensity proxy, subject to shared capacity, buffering, commissioning, and post-construction sustainment constraints.
		
		Rather than seeking a universally optimal strategy, the analysis identifies the operating regimes in which different architectures become competitive. Under the adopted baseline stress-test parameterization, direct rocket delivery completes construction in approximately 8.01 years, whereas the purely space-elevator-assisted architecture requires approximately 235.91 years. At the same fast-build benchmark, the hybrid solution remains effectively rocket-dominant, and disruption tests show that rocket-side rather than SE-side limitations govern the baseline neighborhood. More generally, movement toward SE-assisted operation occurs only when schedule pressure is relaxed and Earth-launch burden is weighted more strongly. Reduced-model diagnostics further indicate no finite-horizon Earth--Moon launch phase-out in the adopted baseline neighborhood or within the explored reduced domain.
		
		These results recast lunar logistics architecture selection as a regime-identification problem rather than a search for a universally superior transport mode. The proposed framework provides a quantitative basis for evaluating when space-elevator-assisted transport is negligible, when it becomes more relevant along the trade-off set, and when no-phaseout persistence remains the dominant reduced-model outcome.
		
		\noindent\textbf{Keywords:} lunar logistics; space elevator; logistics architecture; reliability analysis; multi-objective optimization; transshipment buffering
	\end{abstract}
	
	\par\medskip
	
	\section{Introduction}
	
	For clarity, the abbreviations used in this section are Apex (orbital transshipment node), CTMC (continuous-time Markov chain), CI (cost index), LI (launch-intensity index), and SE (space elevator).
	
	Future lunar infrastructure programs are likely to depend on sustained logistics systems capable of transporting very large material volumes from Earth to the lunar surface over extended planning horizons. In such settings, system performance is governed not only by nominal transport capacity, but also by the temporal coordination of multiple transport legs, interruptions in service, and limits on intermediate storage. The central question is therefore not simply how to optimize a fixed transport chain, but under what conditions alternative logistics architectures become competitive at the campaign scale \cite{COMAP2026B,Merancy2025ADD,Azami2024LBMC,Xia2025LunarBaseInfra,Vasudeva2026LunarConstruction}.
	
	Among the candidate architectures for high-mass lunar transport, SE-assisted logistics is of particular interest because it could, in principle, reduce reliance on repeated direct Earth--Moon launches \cite{Smitherman2000SpaceElevators,Edwards2000NIAC,Chang2022,Nixon2023TetherMaterials,Swan2023ApexAnchorInitialResearch,Swan2025GalacticHarbours}. However, the competitiveness of such architectures cannot be inferred from nominal lift capacity alone. Introducing an orbital transshipment node may increase routing flexibility, but it also creates intermediate inventory states whose admissibility depends on finite buffer capacity, synchronization between inbound and outbound flows, and the persistence of service on both transport legs. Consequently, even when average annual capacity appears sufficient, a transfer-assisted architecture may fail to deliver competitive system-level performance once storage and timing constraints are taken into account.
	
	A further source of uncertainty is the operational reliability of the SE system. Downtime, degradation, and partial-capacity states may reduce effective throughput enough to alter both feasibility and competitiveness within a prescribed construction horizon. This issue is particularly important for campaign-scale planning, where long-run architecture performance depends not only on nominal design capacity but also on the cumulative effect of reliability-mediated throughput loss. In the present study, SE operation is represented by a small-state CTMC reliability model, and the resulting steady-state availability is introduced into the logistics framework through a reduced-form effective-throughput multiplier. This treatment is intended for architecture-level comparison over long horizons rather than short-timescale dispatch or outage scheduling \cite{Norris1998MarkovChains}.
	
	Existing discussions of lunar transport architectures often emphasize nominal capacity, direct feasibility, or static comparisons of alternative pathways. By contrast, less attention has been paid to the coupled effects of transshipment buffering, reliability degradation, commissioning delay, and schedule pressure on architecture competitiveness. In parallel, the broader space-logistics literature has increasingly moved toward campaign-level network and architecture optimization, including multicommodity flow formulations, dynamic campaign design, and recent review and Earth--Moon logistics frameworks \cite{Ishimatsu2016EMMLogistics,Ho2016FlexiblePath,Ho2024Review,Apa2025EarthMoonLogistics}. This gap motivates a regime-oriented perspective: instead of asking which architecture is universally best, the more informative question is when a given architecture becomes competitive and under what constraints that competitiveness can be sustained.
	
	To address this gap, we consider a representative post-2050 stress-test scenario with a cumulative delivered-mass target of $10^8$ metric tons to the lunar surface \cite{COMAP2026B}. This setting is used as a methodological benchmark rather than as a forecast of any specific mission architecture. Within a unified formulation, we compare three classes of strategies under a common set of feasibility constraints: direct rocket delivery from Earth to the Moon, a two-stage SE-assisted pathway from Earth to Apex and from Apex to the Moon, and hybrid strategies that distribute flow across both pathways over time. In this setting, the Apex node induces a finite-buffer transshipment problem, so that architecture performance depends jointly on throughput, inventory admissibility, commissioning delay, and schedule tolerance.
	
	Because architecture choice involves competing planning priorities, the problem is formulated as a multi-objective optimization task rather than reduced to a single subjective scalarization. Non-dominated trade-offs are evaluated among construction completion time, a normalized cost index, and an Earth-launch-intensity index that serves as a convex proxy for sustained Earth-launch burden. This proxy is motivated by the growing literature showing that launch and re-entry activity carries non-negligible atmospheric and climate implications and that these impacts may intensify as launch cadence increases \cite{Ryan2022,Brown2023SustainableLaunch,Sirieys2022AtmosphericImpact,Barker2024Inventory,Sharma2024Atmosphere,Revell2025Ozone}. To ensure that construction-feasible strategies remain operationally meaningful beyond the build phase, the framework also imposes a one-year post-construction water-resupply requirement for a $100{,}000$-person lunar colony, consistent with NASA life-support and water-management references \cite{Ewert2022BVAD,NASA2023WaterHumanConsumptionTB,NASA2023ECLSSBrief,Williamson2023ISSWater,NASASTD3001V2RevE2025}.
	
	Accordingly, this study addresses three linked questions: under what combinations of throughput, reliability, buffering capacity, and schedule flexibility SE-assisted logistics can become competitive; whether hybridization fundamentally changes the governing bottleneck of the campaign-scale system; and whether reduced-model regime-oriented diagnostics indicate a finite-horizon transition or instead support no-phaseout persistence within the adopted baseline neighborhood and explored reduced domain. To answer these questions, the analysis is organized in layers. A campaign-scale optimization model is used for direct comparison of architecture classes under shared feasibility constraints. Campaign-scale sensitivity and numerical-consistency checks are then used to assess local parameter influence and reproducibility of the principal comparisons. Finally, reduced-model diagnostics are introduced to probe whether finite-horizon Earth--Moon launch phase-out emerges under mature-regime perturbations and reduced-domain scans; these reduced diagnostics are interpreted as mechanism-oriented regime probes rather than as substitutes for full campaign-scale comparisons.
	
	Against this background, the paper makes three contributions. First, it develops a unified flow-and-buffer formulation for direct, transshipment, and hybrid lunar logistics under shared capacity, buffering, commissioning, and sustainment constraints. Second, it incorporates SE reliability through an availability-scaled effective-throughput treatment and quantifies how reliability modifies trade-offs among completion time, cost, and Earth-launch burden, while supporting the principal campaign-scale comparisons with local sensitivity and numerical-consistency checks. Third, it reframes lunar logistics architecture comparison as a regime-identification problem, identifies a rocket-dominant baseline fast-build regime under the adopted parameterization, and uses reduced-model diagnostics to test for finite-horizon transition behavior while documenting no-phaseout persistence in the explored reduced domain.

\section{Methods}

The main abbreviations used in this section are SE (space elevator), CTMC (continuous-time Markov chain), Apex (orbital transshipment node), CI (normalized cost index), and LI (launch-intensity index).

\subsection{Planning scope, transport classes, and parameter provenance}
\label{sec:methods_problem}

We study campaign-scale lunar construction logistics beginning in year 2050 (set as $t=0$) with cumulative construction target
\[
M_{\mathrm{goal}} = 10^8 \ \text{t}.
\]
The baseline planning envelope follows the stress-test setting in \cite{COMAP2026B}. Throughout, however, this envelope is used as a methodological benchmark for architecture comparison rather than as a forecast of a specific post-2050 mission design. All headline architecture-comparison claims reported in the main Results section are drawn from the campaign-scale model; reduced-model results are used only for mechanism-oriented diagnostics.

To improve interpretability and reproducibility, all model inputs are partitioned into three provenance classes: \emph{scenario-defined} quantities adopted directly from the stress-test setting, \emph{literature-informed} quantities cross-checked against open technical sources, and \emph{calibration/decision} quantities introduced to represent campaign-scale planning constraints or within-study comparative surrogates. This partition is used to distinguish externally specified assumptions from quantities introduced for controlled architecture comparison. Unless explicitly stated otherwise, the sensitivity ranges reported for calibration/decision quantities are exploratory perturbation ranges used for robustness checks; they are not interpreted as statistical confidence intervals, posterior intervals, or probabilistic priors.

Within a common planning framework, we compare three transport architecture classes:
(i) direct rocket delivery from Earth to the Moon,
(ii) SE-assisted two-stage delivery from Earth to the Apex and then from the Apex to the Moon, and
(iii) hybrid policies that distribute flow across both routes over time.
The reference infrastructure contains three identical ports. Because the present objective is campaign-scale architecture comparison rather than port-level dispatch, flows and inventories are aggregated into system totals. Under this aggregate formulation, the state variables are the cumulative mass delivered to the lunar surface, $M(t)$ (t), and the total inventory stored at the Apex, $A(t)$ (t). The Apex is modeled as a finite transshipment buffer with capacity $A_{\max}$.

The admissible controls are the Earth$\to$Moon launch rate $F_{EM}(t)$ (launches/year), the Apex$\to$Moon launch rate $F_{AM}(t)$ (launches/year), and the Earth$\to$Apex elevator flow $\Phi_{SE}(t)$ (t/year), subject to
\begin{equation}
	0 \leq F_{EM}(t) \leq F_{EM}^{\max}, \qquad
	0 \leq F_{AM}(t) \leq F_{AM}^{\max}, \qquad
	0 \leq \Phi_{SE}(t) \leq \Phi_{SE}^{\max}(t).
	\label{eq:controls_bounds_m3}
\end{equation}
Each rocket launch delivers payload mass $L_R$ (t/launch), yielding the rocket mass flows
\begin{equation}
	\Phi_{EM}(t)=L_R\,F_{EM}(t), \qquad
	\Phi_{AM}(t)=L_R\,F_{AM}(t).
	\label{eq:rocket_flows_methods}
\end{equation}
The scenario specifies $L_R\in[100,150]$ with nominal $L_R=125$ \cite{COMAP2026B}. Open launch-system documentation is used only for order-of-magnitude plausibility checks rather than for exact architectural matching \cite{SpaceX2025FalconUsersGuide,NASA2023SLSFactsheet}. Throughout, mass is measured in t, mass flow in t/year, and launch rates in launches/year. Any fixed rescalings used in tables or figures are presentational only.

Table~\ref{tab:params_methods} summarizes the baseline inputs used in the main simulations and classifies each quantity by provenance. Parameters labeled as calibration/decision quantities should be interpreted as within-study design or policy surrogates. Their role is not to provide absolute techno-economic prediction, but to support comparative regime analysis under transparent and perturbable assumptions.

\begin{table}[htbp]
	\centering
	\caption{Baseline inputs used in the main simulations and their provenance classes. Sensitivity ranges denote exploratory perturbation ranges for robustness checks, not statistical uncertainty intervals.}
	\label{tab:params_methods}
	\setlength{\tabcolsep}{3pt}
	\renewcommand{\arraystretch}{1.0}
	\scriptsize
	\resizebox{\linewidth}{!}{%
		\begin{tabular}{@{}l p{2.4cm} c c p{4.2cm}@{}}
			\toprule
			Symbol & Meaning & Baseline & Sensitivity range & Provenance and role \\
			\midrule
			$M_{\mathrm{goal}}$ & Construction mass target (t) & $10^8$ & -- &
			Scenario-defined stress-test quantity \cite{COMAP2026B} \\
			$P$ & Operational population & $100{,}000$ & -- &
			Scenario-defined sustainment quantity \cite{COMAP2026B} \\
			$K_{\mathrm{SE}}^{\mathrm{sys}}$ & SE throughput (system total, t/yr) & $537{,}000$ & alt.\ $179{,}000$ &
			Scenario-defined value with an explicit system-interpretation ambiguity; the baseline adopts the per-port reading and the whole-system reading is evaluated as a dedicated alternative case \cite{COMAP2026B} \\
			$\alpha$ & SE ramp-up initial fraction & $0.02$ & $[0.018,0.022]$ &
			Calibration quantity for the commissioning profile in Eq.~\eqref{eq:logistic_ramp} \\
			$r$ & SE ramp-up growth rate (yr$^{-1}$) & $0.18$ & $[0.162,0.198]$ &
			Calibration quantity for the commissioning profile in Eq.~\eqref{eq:logistic_ramp} \\
			$\bar{\eta}$ & Mean effective availability & $0.8714$ & $[0.7843,0.9585]$ &
			Planning-scale reduced-form reliability multiplier derived from an auxiliary CTMC interpretation \\
			$L_R$ & Rocket payload to Moon (t/launch) & $125$ & $[100,150]$ &
			Scenario input with literature-based plausibility checks \cite{SpaceX2025FalconUsersGuide,NASA2023SLSFactsheet} \\
			$A_{\max}$ & Apex buffer capacity (t) & $5.0\times 10^5$ & $[4.5,5.5]\times 10^5$ &
			Calibration/decision parameter representing finite transshipment storage \\
			$F_{EM}^{\max}$ & Max Earth$\to$Moon cadence (launch/yr) & $1.0\times 10^5$ & $[0.9,1.1]\times 10^5$ &
			Calibration/decision cap representing campaign-scale launch-policy or operations limits \\
			$F_{AM}^{\max}$ & Max Apex$\to$Moon cadence (launch/yr) & $2.0\times 10^5$ & $[1.8,2.2]\times 10^5$ &
			Calibration/decision cap set sufficiently above the baseline downstream use level to avoid trivially dominating the baseline comparison \\
			$w_{\mathrm{gross}}$ & Gross water use (kg/(person$\cdot$day)) & $3.75$ & $[3.0,4.5]$ &
			Literature-informed bound \cite{NASA2023WaterHumanConsumptionTB,NASASTD3001V2RevE2025,NASA2023ECLSSBrief,Williamson2023ISSWater} \\
			$\rho$ & Water recovery fraction & $0.95$ & $[0.90,0.98]$ &
			Literature-informed bound \cite{NASA2023ECLSSBrief,Williamson2023ISSWater} \\
			$(C_{EM},C_{AM},c_{SE})$ & Relative unit-cost coefficients & $(1.0\!\times\!10^8,\ 2.0\!\times\!10^7,\ 250)$ & $\pm 10\%$ each &
			Calibration coefficients defining a within-study comparative cost surrogate; they are not interpreted as absolute techno-economic forecasts \\
			$\Delta t$ & Time step (yr) & $1/12$ & -- &
			Monthly planning grid \\
			\bottomrule
		\end{tabular}%
	}
\end{table}

\subsection{Campaign-scale flow dynamics, commissioning, and reduced-form reliability}
\label{sec:methods_dynamics}

The Apex inventory is the coupling state between upstream elevator inflow and downstream transshipment launches. The aggregate campaign-scale dynamics are
\begin{align}
	\dot A(t) &= \Phi_{SE}(t)-\Phi_{AM}(t), \qquad 0 \leq A(t) \leq A_{\max},
	\label{eq:buffer_methods}\\
	\dot M(t) &= \Phi_{EM}(t)+\Phi_{AM}(t), \qquad M(0)=0.
	\label{eq:mass_methods}
\end{align}
Equation~\eqref{eq:buffer_methods} couples lift, storage, and downstream release: buffer saturation limits upstream inflow, whereas buffer depletion limits Apex$\to$Moon launch activity.

SE commissioning is modeled by a bounded monotone ramp
\begin{equation}
	N_{\mathrm{SE}}(t)=\frac{1}{1+\left(\frac{1-\alpha}{\alpha}\right)e^{-rt}},
	\label{eq:logistic_ramp}
\end{equation}
with baseline $(\alpha,r)=(0.02,\,0.18~\mathrm{yr}^{-1})$. This two-parameter logistic form is used because it provides a smooth commissioning profile with an interpretable early-stage fraction $\alpha$ and characteristic growth timescale $r^{-1}$ while avoiding additional discontinuity parameters.

The scenario statement reports a perfect-condition lift figure of $179{,}000$ t/year \cite{COMAP2026B}. Because three identical ports are specified but only one lift value is reported, the baseline interpretation treats $179{,}000$ t/year as a per-port figure, yielding
\begin{equation}
	K_{\mathrm{SE}}^{\mathrm{sys}} = 3\times 179{,}000\ \text{t/year}.
	\label{eq:Ksys_methods}
\end{equation}
The whole-system interpretation is examined separately as an alternative scenario rather than merged into the baseline.

Reliability is incorporated at the planning scale through a reduced-form effective-availability factor. Let $X(t)$ be an irreducible finite-state CTMC with state space $\{1,\dots,n\}$, generator matrix $Q$, and state-dependent capacity multipliers $\gamma_i\in[0,1]$. If $\pi$ is the stationary distribution satisfying $\pi Q=0$ and $\sum_i \pi_i=1$, then the stationary mean availability is
\begin{equation}
	\bar{\eta}:=\sum_{i=1}^{n}\pi_i\gamma_i.
	\label{eq:eta_bar_methods}
\end{equation}
At the optimization level, only the reduced multiplier $\bar{\eta}$ is propagated into the campaign-scale model, because the objective here is long-horizon architecture comparison rather than short-timescale outage scheduling. The complete CTMC state definitions, transition structure, generator matrix, and capacity multipliers used to compute the baseline $\bar{\eta}$ are reported in the Supplementary Information and archived code package. In the baseline runs,
\[
\bar{\eta}=0.8714,
\]
which implies
\[
K_{\mathrm{SE}}^{\mathrm{eff}}
=\bar{\eta}\,K_{\mathrm{SE}}^{\mathrm{sys}}
=467{,}921\ \text{t/year}.
\]

For scenario-level disruption cases, temporary degradations are represented by smooth multiplicative health windows. For transport mode $q\in\{\mathrm{SE},\mathrm{R}\}$, we write
\begin{equation}
	h_q(t)=1-s_q\Big[\sigma_k(t-t_{0,q})-\sigma_k\big(t-(t_{0,q}+d_q)\big)\Big],
	\label{eq:health_window_methods}
\end{equation}
where $s_q\in[0,1]$ is the disruption severity, $t_{0,q}$ is the start time, $d_q$ is the window duration, and $\sigma_k$ is the smooth logistic step. In the main reported disruption scenarios, failure timing is specified in relative form so that start times and durations are defined as fractions of the relevant case-horizon seed; this preserves comparable disruption placement across cases with different schedule lengths.

Accordingly, the admissible elevator flow is constrained by
\begin{equation}
	0\le \Phi_{SE}(t)\le h_{\mathrm{SE}}(t)\,N_{\mathrm{SE}}(t)\,\bar{\eta}\,K_{\mathrm{SE}}^{\mathrm{sys}},
	\label{eq:se_effective_methods}
\end{equation}
with $h_{\mathrm{SE}}(t)\equiv 1$ in non-disruption cases. Analogously, rocket-side disruption modifies the effective launch cap through $F_{EM}^{\max}\mapsto h_{\mathrm{R}}(t)\,F_{EM}^{\max}$, with $h_{\mathrm{R}}(t)\equiv 1$ when no rocket-side failure is imposed.

The construction completion time in continuous time is the first hitting time
\begin{equation}
	M(T_{\mathrm{build}})=M_{\mathrm{goal}}.
	\label{eq:build_target_methods}
\end{equation}
Because direct optimization over this hitting time is non-smooth, the discretized problem instead enforces completion by a prescribed deadline and recovers the achieved completion time from the optimized trajectory.

To connect construction feasibility with near-term sustainment, we additionally impose a one-year post-construction water-resupply requirement \cite{COMAP2026B}. Let $P = 100{,}000$ denote the colony population. Gross daily water use is modeled as
\begin{equation}
	w_{\mathrm{gross}} \in [3.0, 4.5]\ \mathrm{kg/(person\cdot day)},
	\label{eq:wgross_methods}
\end{equation}
and the corresponding net annual resupply mass under recovery fraction $\rho \in [0.90, 0.98]$ is
\begin{equation}
	W_{\mathrm{net}}
	= (1-\rho)\, P\, w_{\mathrm{gross}} \frac{365}{1000}
	\quad \text{t/year}.
	\label{eq:Wnet_methods}
\end{equation}
This requirement is used as an operational admissibility screen linking the build phase to near-term post-construction sustainment. Additional consumables beyond water are considered only in extended sensitivity cases using traceable bounds \cite{Ewert2022BVAD}.

\subsection{Discretization, optimization tasks, and reporting metrics}
\label{subsec:workflow}

For computation, time is discretized with step $\Delta t$ (years), and controls are assumed piecewise constant on each interval $[k\Delta t,(k+1)\Delta t)$. Let $A_k\approx A(k\Delta t)$ and $M_k\approx M(k\Delta t)$, and let $(F_{EM,k},F_{AM,k},\Phi_{SE,k})$ denote the per-period controls. The discrete dynamics are
\begin{align}
	A_{k+1}&=A_k+\Delta t\left(\Phi_{SE,k}-L_R\,F_{AM,k}\right),\qquad 0\le A_{k+1}\le A_{\max},
	\label{eq:disc_buffer_m3}\\
	M_{k+1}&=M_k+\Delta t\left(L_R\,F_{EM,k}+L_R\,F_{AM,k}\right),
	\label{eq:disc_mass_m3}
\end{align}
for $k=0,\ldots,K-1$.

The computational workflow addresses three distinct tasks:
(i) determination of the minimum feasible build time,
(ii) construction of near-fast-build cost comparators, and
(iii) broader trade-off exploration through frontier sweeps.

Let $K_b=\lceil \bar{T}/\Delta t\rceil$ and $K_w=\lceil 1/\Delta t\rceil$, where $\bar{T}$ denotes a prescribed construction deadline. Completion by the deadline is enforced by
\begin{equation}
	M_{K_b}\ge M_{\mathrm{goal}}.
	\label{eq:completion_by_deadline_m3}
\end{equation}
The reported grid-based completion time is
\[
T_{\mathrm{rep}}:=\min\{k\Delta t:\ M_k\ge M_{\mathrm{goal}}\}\le \bar{T}.
\]
The one-year post-construction water-resupply requirement is enforced over the subsequent window,
\begin{equation}
	M_{K_b+K_w}-M_{K_b}\ge W_{\mathrm{net}},
	\label{eq:water_after_deadline_m3}
\end{equation}
which guarantees completion no later than $K_b\Delta t$ and sustains a full year of post-build water demand under the optimized trajectory.

Three admissible transport architecture classes are imposed through control restrictions. In the rocket-only case, $\Phi_{SE,k}=0$ and $F_{AM,k}=0$ for all $k$. In the SE-assisted two-stage case, $F_{EM,k}=0$ for all $k$. In the hybrid case, only the shared feasibility constraints apply. Under identical inputs, the hybrid class contains the two restricted classes as special cases.

For each feasible trajectory, we report the build completion time $T_{\mathrm{rep}}$ (years), the total Earth-launch count
\[
N_E:=\sum_{k=0}^{K_b+K_w-1}\Delta t\,F_{EM,k},
\]
the raw Earth-launch intensity
\[
J_E:=\sum_{k=0}^{K_b+K_w-1}\Delta t\,F_{EM,k}^2,
\]
and the raw trajectory cost
\[
J_C^{\mathrm{raw}}
:=\sum_{k=0}^{K_b+K_w-1}\Delta t\left(C_{EM}F_{EM,k}+C_{AM}F_{AM,k}+c_{SE}\Phi_{SE,k}\right).
\]
Here $J_E$ is used as a convex operational proxy for sustained Earth-launch burden; it is not interpreted as a direct emissions inventory, atmospheric chemistry model, or life-cycle environmental assessment. For within-study cross-scenario reporting, we define
\[
\mathrm{LI}:=\frac{J_E}{J_E^{\mathrm{ref}}},
\qquad
\mathrm{CI}:=\frac{J_C^{\mathrm{raw}}}{J_C^{\mathrm{ref}}},
\]
where $J_E^{\mathrm{ref}}>0$ and $J_C^{\mathrm{ref}}>0$ are fixed reference values within a stated comparison family. The reference pair used for each figure or table is explicitly identified in the corresponding caption or archived metadata. Accordingly, $\mathrm{LI}$ and $\mathrm{CI}$ are normalized reporting indices for architecture comparison, whereas the solver operates on dimensionally scaled but order-equivalent implementation objectives to improve numerical conditioning.

Because $M_{\mathrm{goal}}/L_R$ implies on the order of $10^6$ launches, relaxing discrete launch counts to bounded launch-rate controls is accurate at the campaign scale. Under nearest-integer rounding, the per-period count error is at most $1/2$ launches per leg. Hence the absolute cumulative delivery error is bounded by $\tfrac12 K L_R$ over $K$ periods, and the corresponding relative error is bounded by $\tfrac12 K L_R/M_{\mathrm{goal}}$, which is negligible at the scale considered here.

For baseline architecture comparison and representative scenario reporting, the first task is to compute the minimum feasible build-time benchmark $\tau_\star$ by bisection on the construction horizon. The second task is to construct a near-fast-build cost-oriented comparator by minimizing implementation cost subject to
\[
\tau \le (1+\epsilon_\tau)\tau_\star,\qquad \epsilon_\tau=0.05.
\]
The third task is to examine broader cost--launch-burden trade-offs through frontier sweeps described below.

\subsection{Frontier construction and regime-oriented diagnostics}
\label{sec:methods_frontier}

To characterize broader architecture-dependent trade-offs in the campaign-scale model, we use an $\varepsilon$-constraint construction. Specifically, implementation cost is minimized subject to a launch-intensity bound,
\[
\min J_C^{\mathrm{raw}}
\quad
\text{s.t.}
\quad
J_E \le \varepsilon ,
\]
with $\varepsilon$ swept over a predefined grid. The resulting objective vectors are filtered using Pareto non-dominance. A solution $\mathbf{J}^{(1)}$ dominates $\mathbf{J}^{(2)}$ if $J_i^{(1)}\le J_i^{(2)}$ for all $i\in\{T,\mathrm{CI},\mathrm{LI}\}$ and strict inequality holds for at least one component. A small numerical tolerance is used in post-processing to avoid spurious dominance induced by solver tolerances; the exact tolerance used in each reported sweep is fixed in the archived scripts and metadata.

To formalize the architecture labels used in the results, define the cumulative lunar-delivered masses over the reporting horizon
\[
S_R:=\sum_{k=0}^{K_b+K_w-1}\Delta t\,L_R F_{EM,k},
\qquad
S_{SE}:=\sum_{k=0}^{K_b+K_w-1}\Delta t\,L_R F_{AM,k}.
\]
The route-share statistic is then
\[
\chi_R:=\frac{S_R}{S_R+S_{SE}}.
\]
A trajectory is classified as \emph{rocket-dominant} if $\chi_R\ge 0.9$, \emph{SE-assisted} if $\chi_R\le 0.1$, and \emph{mixed} otherwise. These thresholds are used only for regime labeling and visualization; they do not alter the optimization problem itself.

In addition to the campaign-scale formulation, we use a reduced mature-regime diagnostic model for mechanism-oriented analysis. This reduced model retains the flow-balance, capacity, and route-allocation structure of the campaign-scale formulation, but removes early commissioning transients and finite disruption windows so that mature-regime throughput--launch-aversion interactions can be examined in isolation. In the archived reduced-model diagnostics, the mature-regime baseline sets $\alpha=1$, deactivates finite disruption windows, and uses a fixed 50-year horizon. The reduced Monte Carlo diagnostic perturbs \(\{K_{\max,\mathrm{total}},L_R,r,C_E,C_A\}\) independently by \(\pm5\%\) around the adopted reduced baseline, with requested launch-aversion parameter \(\gamma=0.08\) and terminal-mass screening based on \(M_T\ge 0.99\,M_{\mathrm{goal}}\). The reduced 2D scan sweeps effective-throughput factors and launch-aversion multipliers over a predefined grid around a reference reduced-model Earth-launch-burden level. These reduced diagnostics are not used to generate the primary architecture-comparison claims in this paper; rather, they are used only to test whether finite-horizon launch phase-out emerges or whether no-phaseout persistence remains dominant within the explored reduced domain.

\subsection{Numerical solution and verification protocol}
\label{sec:methods_numerics}

The computational workflow uses two linked model layers: (i) a campaign-scale full model for the main architecture comparisons, representative policy variants, and local sensitivity analyses; and (ii) a reduced mature-regime diagnostic model for launch phase-out and regime-oriented sensitivity analysis. These two layers share data-processing conventions and archive formats, but they are not numerically identical and are interpreted separately in the results.

The discretized nonlinear programs are solved by direct transcription using CasADi with IPOPT. Horizon search for $\tau_\star$, launch-intensity sweeps, and continuation across neighboring parameter cases are implemented at the script level. Exact solver options, stopping tolerances, bisection tolerances, initialization rules, and restart settings are fixed across each experiment family and documented in the archived code package.

Verification includes four complementary checks. First, randomized-initialization reruns are used to assess solver stability on representative campaign-scale instances. Second, grid-refinement checks compare the monthly baseline grid against finer discretizations on selected cases. Third, limiting-case checks verify expected behavior in the no-SE limit, under relaxed schedule constraints, and under enlarged buffer capacity. Fourth, reduced-form reliability stress tests compare cases with matched $\bar{\eta}$ but different outage timing structures to assess whether campaign-level conclusions depend materially on failure clustering rather than on mean effective capacity alone.

A numerical check is considered passed when (i) the feasibility classification is unchanged and (ii) the corresponding reported comparison metrics remain within predefined tolerances under the relevant perturbation. In the archived implementation, solver outcomes are classified separately as accepted, infeasible, or numerically unsuccessful according to fixed acceptance rules based on solver status, feasibility residuals, and problem-specific admissibility checks. For the reduced Monte Carlo diagnostic, a sample is counted as numerically accepted only when its solver class is acceptable and its terminal delivered mass satisfies \(M_T \ge 0.99\,M_{\mathrm{goal}}\). The exact rules and thresholds used in the automated audits are reported in the Supplementary Information and archived metadata.

Complete machine-readable scenario summaries, representative time-series trajectories, randomized-restart audit logs, and reduced-model robustness outputs are described in the Data Availability and Code Availability statements.


\section{Scenario Results, Trade-offs, and Regime Behavior}
\label{sec:results}

We compare three transport architecture classes under shared baseline assumptions on capacity, buffering, commissioning, and reliability: direct rocket delivery from Earth to the Moon, SE-assisted two-stage delivery from Earth to the Apex and then to the Moon, and hybrid time-varying logistics \cite{COMAP2026B}. This section reports the campaign-scale results that support the paper’s primary comparative claims. Unless stated otherwise, all reported metrics, feasibility conditions, and optimization settings follow Sections~\ref{sec:methods_dynamics}--\ref{sec:methods_numerics}. All headline results in this section are obtained from the campaign-scale model and should be interpreted within the adopted stress-test parameterization rather than as forecasts of a specific mission campaign. Reduced-model diagnostics are considered later only as mechanism-oriented complements for testing transition behavior; they are not used to define the primary comparative claims reported here. The representative campaign-scale benchmarks discussed here are interpreted together with separate numerical-consistency audits, which showed stable repeated-solve behavior for the mixed baseline case under perturbed initializations.

\subsection{Baseline comparison at the minimum-feasible build horizon}
\label{subsec:baseline_tau}

Table~\ref{tab:baseline_strategy_tau} compares the three architecture classes at the minimum-feasible build horizon $\tau_\star$, obtained by bisection on the construction deadline. This benchmark isolates the shortest feasible schedule attainable within each strategy class under the shared baseline assumptions and therefore provides the clearest campaign-scale view of the active fast-build bottleneck.

\begin{table}[htbp]
	\centering
	\caption{Baseline architecture comparison at the minimum-feasible build horizon $\tau_\star$.
		Launch totals are integrated on $[0,T_b+1]$.
		SE throughput [Mt] denotes cumulative Earth$\to$Apex elevator lift over $[0,T_b+1]$
		(not lunar-delivered mass). Cost is reported as total integrated trajectory cost in the implementation units.}
	\label{tab:baseline_strategy_tau}
	
	\small
	\setlength{\tabcolsep}{3pt}
	\renewcommand{\arraystretch}{1.1}
	
	\resizebox{\linewidth}{!}{%
		\begin{tabular}{lcccccc}
			\hline
			Strategy & $T_b$ [yr] & Cost [$10^9$]
			& Earth$\to$Moon launches
			& Apex$\to$Moon launches
			& SE throughput [Mt]
			& $A_{\mathrm{peak}}$ [kt] \\
			\hline
			Rocket-only & 8.01 & 79878.22 & 798782.22 & 0.00    & 0.000 & 0.0 \\
			SE-assisted & 235.91 & 16026.31 & 0.00 & 800065.20 & 100.006 & 8.8 \\
			Hybrid      & 8.01 & 79769.16 & 797407.31 & 1419.28 & 0.169 & 156.9 \\
			\hline
	\end{tabular}}
\end{table}

Three campaign-scale features are apparent from Table~\ref{tab:baseline_strategy_tau}. First, the rocket-only architecture attains the shortest feasible schedule, with $T_b=8.01$ years. This value is close to the direct-capacity lower bound
\[
\frac{M_{\mathrm{goal}}}{L_R F_{EM}^{\max}}
=
\frac{10^8}{125\times 10^5}
=8.0\ \text{years},
\]
which is consistent with a fast-build regime governed primarily by the Earth$\to$Moon launch-capacity limit. Second, the SE-assisted architecture requires $T_b=235.91$ years, indicating that the baseline two-stage route is much slower under the adopted throughput and staging conditions. The simple effective-throughput estimate
\[
\frac{M_{\mathrm{goal}}}{K_{\mathrm{SE}}^{\mathrm{eff}}}
=
\frac{10^8}{467{,}921}
\approx 213.7\ \text{years}
\]
captures the relevant order of magnitude, while the optimized campaign-scale solution is longer because the full formulation additionally enforces staging, finite-buffer, discretization, and post-construction admissibility conditions. Third, the hybrid solution at $\tau_\star$ is effectively indistinguishable from the rocket-only case in build time, while exhibiting only limited SE participation relative to the overall campaign scale. Under the route-share definition in Section~\ref{sec:methods_frontier}, this solution is therefore classified as rocket-dominant.

Table~\ref{tab:baseline_diag_tau} reports derived cadence and effective-throughput diagnostics inferred from Table~\ref{tab:baseline_strategy_tau}. These quantities are not additional optimization outputs; rather, they summarize the implied operating levels associated with the baseline solutions and make the active bottleneck at the minimum-feasible build horizon more explicit.

\begin{table}[htbp]
	\centering
	\caption{Cadence and effective-throughput diagnostics inferred from Table~\ref{tab:baseline_strategy_tau}.
		We report mean cadences $\overline{F}_{EM}=N_{EM}/(T_b+1)$ and $\overline{F}_{AM}=N_{AM}/(T_b+1)$ (launches/year).
		The implied SE throughput is $K_{SE}^{\mathrm{imp}}=(10^6\,M_{SE})/(T_b+1)$ in t/year (with $M_{SE}$ in Mt),
		and the implied effective-throughput ratio is $\delta^{\mathrm{imp}}=K_{SE}^{\mathrm{imp}}/K_{\mathrm{SE}}^{\mathrm{sys}}$.
		Here $K_{\mathrm{SE}}^{\mathrm{sys}}=537{,}000$ t/year is the baseline nominal system throughput
		(Table~\ref{tab:params_methods}).}
	\label{tab:baseline_diag_tau}
	
	\small
	\setlength{\tabcolsep}{4pt}
	\renewcommand{\arraystretch}{1.08}
	
	\begin{tabular}{lrrrr}
		\hline
		Strategy & $\overline{F}_{EM}$ [launch/yr] & $\overline{F}_{AM}$ [launch/yr] & $K_{SE}^{\mathrm{imp}}$ [t/yr] & $\delta^{\mathrm{imp}}$ [-] \\
		\hline
		Rocket-only & 88603.37 & 0.00    & 0        & 0.0000 \\
		SE-assisted & 0.00     & 3372.88 & 422134   & 0.7861 \\
		Hybrid      & 88450.85 & 157.39  & 18748    & 0.0349 \\
		\hline
	\end{tabular}
\end{table}

The diagnostics support the same campaign-scale interpretation. In the SE-assisted case, the implied throughput is of the same order as the baseline effective capacity, which is consistent with elevator-side effective throughput acting as the active limitation for the two-stage route. In the hybrid $\tau_\star$ solution, by contrast, the implied SE utilization remains small relative to system-scale throughput. At the displayed precision, the baseline minimum-feasible schedule is therefore not a materially mixed operating regime; instead, it is a rocket-dominant solution in which the SE leg appears only at secondary levels.

\subsection{Near-minimum-duration hybrid comparator and long-horizon low-burden endpoint}
\label{subsec:tradeoff_endpoints}

To illustrate the extremes of the baseline hybrid trade-off set without reproducing the full frontier, Table~\ref{tab:mixed_tradeoff} reports two representative hybrid endpoints: a near-minimum-duration cost-oriented comparator and a long-horizon low-burden endpoint. The former minimizes implementation cost subject to the near-minimum-duration condition introduced in Section~\ref{subsec:workflow}, whereas the latter suppresses direct Earth-launch use and accepts a much longer build horizon. In the reported long-horizon case, the upper schedule bound $T_{\max}=400$ years is active, so that endpoint should be interpreted as a horizon-capped frontier point rather than as an unconstrained optimum.

\begin{table}[htbp]
	\centering
	\caption{Representative hybrid endpoints under baseline conditions.
		The near-minimum-duration comparator minimizes cost subject to $\tau \le 1.05\,\tau_\star$, whereas the long-horizon low-burden endpoint is reported at the active upper horizon bound.
		All totals are computed on $[0,T_b+1]$.}
	\label{tab:mixed_tradeoff}
	
	\small
	\setlength{\tabcolsep}{3pt}
	\renewcommand{\arraystretch}{1.1}
	
	\resizebox{\linewidth}{!}{%
		\begin{tabular}{lcccccc}
			\hline
			Endpoint & $T_b$ [yr] & Cost [$10^9$]
			& Earth$\to$Moon launches
			& Apex$\to$Moon launches
			& SE throughput [Mt]
			& $A_{\mathrm{peak}}$ [kt] \\
			\hline
			Near-minimum-duration cost comparator & 8.41 & 79757.75 & 797269.89 & 1535.73 & 0.185 & 132.2 \\
			Long-horizon low-burden ($T_{\max}$ active) & 400.00 & 16568.48 & 6822.89 & 793069.76 & 99.204 & 333.3 \\
			\hline
	\end{tabular}}
\end{table}

The two endpoints indicate that, under the adopted baseline parameterization, the hybrid class spans the two limiting operating modes rather than producing a distinct regime that is simultaneously fast-building and low-burden. The near-minimum-duration comparator is practically indistinguishable from the rocket-dominant benchmark in build time and still relies overwhelmingly on direct Earth$\to$Moon launches. By contrast, the long-horizon low-burden endpoint sharply suppresses direct Earth-launch use and shifts transport onto the SE-assisted chain, but only at the cost of a much longer construction schedule. Within the present parameterization, the hybrid class therefore acts mainly as a connector between the rocket-dominant and SE-assisted limits as objective emphasis changes, rather than creating a third fast-build, low-burden regime.

Figure~\ref{fig:optimal_switching} shows a representative hybrid schedule in physical units. Rocket cadence is plotted in kLaunch/yr, and SE lift is shown as the equivalent cadence $\widetilde F_{SE,\mathrm{eq}}=(\Phi_{SE}/L_R)/10^3$. The figure is intended to assess whether the optimized hybrid policy exhibits substantial endogenous mode switching or remains predominantly Earth-launch driven over time.

\begin{figure}[htbp]
	\centering
	\includegraphics[width=0.48\linewidth]{figure/Fig1_OptimalSwitching.png}
	\caption{Representative hybrid schedule in physical units.
		Rocket cadence is plotted in kLaunch/yr; SE lift is shown as the equivalent cadence
		$\widetilde F_{SE,\mathrm{eq}}=(\Phi_{SE}/L_R)/10^3$.
		The plot highlights the relative timing of rocket and SE contributions within the optimized hybrid schedule.}
	\label{fig:optimal_switching}
\end{figure}

The schedule profile is consistent with Table~\ref{tab:mixed_tradeoff}: near the fast-build benchmark, endogenous mode switching remains limited and the optimized hybrid trajectory is overwhelmingly Earth-launch driven. More substantial modal substitution appears only after schedule pressure is relaxed or Earth-launch burden is constrained more strongly.

\subsection{Disruption sensitivity under non-perfect operation}
\label{subsec:resilience}

We next test whether non-perfect operation changes the active bottleneck identified in the baseline fast-build regime. The baseline architecture comparisons use the mean-availability formulation, whereas the resilience illustrations and stress tests additionally impose deterministic time-window capacity reductions on the effective capacity profile. In the main reported disruption cases, the time windows are specified in relative form so that degradation episodes remain comparable across cases with different build horizons. These cases are stress tests defined within the campaign-scale model; they are not stochastic simulations of random failure histories.

Figure~\ref{fig:resilience_deliverable} illustrates a representative disruption-and-recovery episode. During the disruption window, the effective capacity of the affected mode is reduced, and the optimized disrupted schedule compensates by reallocating effort to the remaining active leg subject to the corresponding launch-rate or throughput caps. In the illustrated case, the recovery response increases rocket cadence relative to the baseline schedule in order to maintain progress toward the delivered-mass target.

\begin{figure}[htbp]
	\centering
	\includegraphics[width=0.48\linewidth]{figure/Fig2_Resilience.png}
	\caption{Representative disruption-and-recovery episode.
		Top: launch-cadence adjustment during and after a disruption window
		(dashed = baseline schedule, solid = disrupted-scenario optimum).
		Bottom: cumulative lunar-delivered mass $M(t)$ compared to $M_{\mathrm{goal}}$.
		Rocket cadence is plotted in kLaunch/yr.}
	\label{fig:resilience_deliverable}
\end{figure}

Table~\ref{tab:mixed_disrupt_tau} reports hybrid outcomes under three disruption modes at the minimum-feasible build horizon: SE disrupted, rockets disrupted, and both disrupted.

\begin{table}[htbp]
	\centering
	\caption{Hybrid outcomes under disruption modes at the minimum-feasible build horizon.
		SE throughput [Mt] denotes cumulative Earth$\to$Apex lift over $[0,T_b+1]$.
		$T_b$ is the reported build-completion time for the disrupted scenario.}
	\label{tab:mixed_disrupt_tau}
	
	\small
	\setlength{\tabcolsep}{3pt}
	\renewcommand{\arraystretch}{1.1}
	
	\resizebox{\linewidth}{!}{%
		\begin{tabular}{lccccc}
			\hline
			Disruption case & $T_b$ [yr] & Cost [$10^9$]
			& Earth$\to$Moon launches
			& SE throughput [Mt]
			& $A_{\mathrm{peak}}$ [kt] \\
			\hline
			Baseline          & 8.01 & 79769.16 & 797407.31 & 0.169 & 156.9 \\
			SE disrupted      & 8.01 & 79780.97 & 797501.36 & 0.174 & 151.9 \\
			Rockets disrupted & 8.45 & 83593.86 & 797079.75 & 0.486 & 181.5 \\
			Both disrupted    & 8.45 & 83619.83 & 797188.13 & 0.487 & 176.4 \\
			\hline
	\end{tabular}}
\end{table}

The near-identical baseline and SE-disrupted rows are informative because, at the minimum-feasible build horizon, the hybrid optimum allocates only limited mass to the SE leg, consistent with Tables~\ref{tab:baseline_strategy_tau} and~\ref{tab:mixed_tradeoff}. Reducing SE availability therefore does not change the reported schedule at the displayed precision. By contrast, rocket-side disruption shifts the minimum-feasible build time from $8.01$ to $8.45$ years and increases total cost, which is consistent with the rocket leg remaining the active fast-build bottleneck in the baseline regime. The combined-disruption case remains close to the rockets-disrupted case for the same reason: when the rocket side is already the active limitation, additional SE degradation has only a secondary effect.

Disruption timing also matters. Table~\ref{tab:early_disrupt} reports early-window stress tests for the near-minimum-duration hybrid comparator of Table~\ref{tab:mixed_tradeoff}. Unlike Table~\ref{tab:mixed_disrupt_tau}, these cases impose an exogenous outage over a fixed early interval and then restore capacity to its baseline pattern.

\begin{table}[htbp]
	\centering
	\caption{Early-window disruption stress test for the near-minimum-duration hybrid comparator.
		SE throughput [Mt] denotes cumulative Earth$\to$Apex lift over $[0,T_b+1]$.}
	\label{tab:early_disrupt}
	
	\small
	\setlength{\tabcolsep}{3pt}
	\renewcommand{\arraystretch}{1.1}
	
	\resizebox{\linewidth}{!}{%
		\begin{tabular}{lcccc}
			\hline
			Early disruption case & $T_b$ [yr] & Cost [$10^9$]
			& Earth$\to$Moon launches
			& SE throughput [Mt] \\
			\hline
			SE disrupted (early)   & 8.41 & 79770.48 & 797409.40 & 0.172 \\
			Both disrupted (early) & 8.87 & 82407.25 & 796851.29 & 0.497 \\
			\hline
	\end{tabular}}
\end{table}

The early-window stress tests support the same mechanism. Early SE disruption alone has little effect because the near-minimum-duration hybrid comparator still uses the SE leg only weakly. By contrast, early joint disruptions increase build time and total cost through rocket-side limitations. The resilience response is therefore regime-dependent: disruption to a weakly used leg has limited consequence, whereas disruption to the active bottleneck shifts both schedule and transport composition.

\subsection{Regime interpretation and transition implications}
\label{subsec:regime_behavior}

Taken together, the baseline comparison, the representative hybrid endpoints, and the disruption stress tests support a consistent campaign-scale regime interpretation. Under the adopted throughput, availability, and cadence assumptions, the fast-feasible region is rocket-dominant. Near the minimum-feasible build horizon, the hybrid class is not materially distinct from rocket-only delivery, and degrading SE performance has little effect because the SE leg is not an active contributor at that operating point.

A separate, conditional conclusion concerns the transition away from this rocket-dominant regime. As schedule pressure is relaxed, or as Earth-launch burden is constrained more strongly, optimized policies move along the feasible trade-off set toward the SE-assisted chain. This shift substantially reduces direct Earth-launch reliance but requires much longer build horizons. The present campaign-scale results therefore do not identify a materially competitive third regime that is simultaneously fast-building and low-burden; rather, within the adopted parameterization, the hybrid class mainly connects the rocket-dominant and SE-assisted limits.

This interpretation is further reinforced by the archived campaign-scale bridge sensitivity study. In that study, the four representative campaign-scale cases \texttt{mixed\_base}, \texttt{mixed\_sefail}, \texttt{mixed\_rocketfail}, and \texttt{mixed\_bothfail} were re-solved under \(\tau_\star\)-oriented and cost-oriented policies while varying \(K_{\max,\mathrm{total}}\), \(\delta_{\mathrm{eff}}\), \(L_R\), \(\gamma_{\mathrm{env}}\), and \(C_E\) by multiplicative factors \(\{0.6,0.8,1.0,1.2,1.5\}\). Across all archived full-model bridge runs, the regime label remained rocket-dominant. The strongest changes were induced by \(L_R\), whereas variations in \(K_{\max,\mathrm{total}}\) and \(\delta_{\mathrm{eff}}\) were comparatively small in the explored campaign-scale grid and did not trigger a regime flip. These bridge results indicate that the baseline fast-build neighborhood is not merely locally rocket-dominant at a single benchmark point, but remains rocket-dominant across a wider set of campaign-scale throughput-, payload-, and cost-side perturbations.

These campaign-scale observations motivate the reduced-model diagnostics discussed later. The reduced layer is used to test whether any finite-horizon transition emerges within the explored reduced domain or whether no-phaseout persistence remains the dominant behavior; it is not used to redefine the campaign-scale conclusions reported in the present section.


\section{Sensitivity, Local Drivers, and Reduced-Model No-Phaseout Diagnostics}
\label{sec:sensitivity}

To assess whether the main findings depend narrowly on the adopted baseline parameterization, we examine how reported performance metrics and operating regimes change under parameter perturbations. The analysis is organized around four complementary diagnostics: campaign-scale one-at-a-time (OAT) sensitivity around a fixed baseline comparator, campaign-scale bridge sensitivity across representative cases and wider parameter factors, reduced-model Monte Carlo diagnostics under joint perturbations, and a reduced mature-regime two-dimensional scan over effective-throughput and Earth-launch-burden conditions. These diagnostics address related but distinct questions. The OAT study quantifies local sensitivity of campaign-scale outputs near the adopted mixed baseline comparator; the campaign-scale bridge study tests whether the same full-model regime interpretation persists under broader key-parameter perturbations across representative baseline and disruption cases; and the reduced-model diagnostics test whether finite-horizon launch phase-out emerges in the adopted reduced baseline neighborhood or across the explored reduced domain. Detailed sampling conventions, phase-out bookkeeping rules, and archived implementation metadata are provided in the Supplementary Information.

\subsection{Campaign-scale local sensitivity and reduced-model joint-perturbation diagnostics}
\label{subsec:tornado_mc}

We first examine local OAT sensitivity around a fixed campaign-scale benchmark and then assess reduced-model behavior under joint perturbations in the mature-regime formulation. For each perturbed parameter \(p\), the local campaign-scale analysis records the relative changes
\[
\Delta_T(p)=\frac{T_b(p)-T_b^{(0)}}{T_b^{(0)}},\qquad
\Delta_{\mathrm{CI}}(p)=\frac{\mathrm{CI}(p)-\mathrm{CI}^{(0)}}{\mathrm{CI}^{(0)}},\qquad
\Delta_E(p)=\frac{J_E(p)-J_E^{(0)}}{J_E^{(0)}},
\]
where \((\cdot)^{(0)}\) denotes the corresponding baseline benchmark value. In the archived implementation, this OAT analysis is conducted around the mixed, failure-free campaign-scale baseline comparator. The tornado plot reports the larger absolute deviation induced by the prescribed positive and negative perturbations. Unless a parameter-specific range is defined in Table~\ref{tab:params_methods}, the OAT study uses the same local perturbation convention as the archived sensitivity batch.

The archived campaign-scale bridge sensitivity study broadens this picture by re-solving the full model across the representative cases \texttt{mixed\_base}, \texttt{mixed\_sefail}, \texttt{mixed\_rocketfail}, and \texttt{mixed\_bothfail}, under both \(\tau_\star\)-oriented and cost-oriented policies, while varying
\[
\{K_{\max,\mathrm{total}},\,\delta_{\mathrm{eff}},\,L_R,\,\gamma_{\mathrm{env}},\,C_E\}
\]
by multiplicative factors \(\{0.6,0.8,1.0,1.2,1.5\}\). This bridge layer is still campaign-scale, but it goes beyond local OAT perturbations by testing whether broader full-model parameter changes can move the system out of the rocket-dominant regime. In the archived bridge grid, all full-model runs remained rocket-dominant. The strongest variations were induced by \(L_R\), while \(C_E\) primarily altered cost and \(K_{\max,\mathrm{total}}\) and \(\delta_{\mathrm{eff}}\) produced comparatively small changes in build time and did not trigger regime reversal within the explored full-model domain.

The reduced Monte Carlo diagnostic is conceptually different. Here selected reduced-model inputs are perturbed simultaneously by independent multiplicative factors around the mature-regime baseline. In the archived run, the perturbed set is
\[
\{K_{\max,\mathrm{total}},\,L_R,\,r,\,C_E,\,C_A\},
\]
with independent \(\pm 5\%\) perturbations, requested launch-aversion parameter \(\gamma=0.08\), and terminal-mass screening based on \(M_T \ge 0.99\,M_{\mathrm{goal}}\). For each sample, the reduced mature-regime problem is re-solved and the resulting feasibility class, Earth-launch burden, and finite-horizon launch phase-out status are recorded. A sample is counted as numerically accepted only when it satisfies the reduced-model acceptance rule summarized in Section~\ref{sec:methods_numerics}. If no phase-out occurs within the sampled horizon, the run is recorded as a no-phaseout case within that horizon rather than assigned an artificial transition time.

Figure~\ref{fig:sens_tornado_mc} combines the campaign-scale OAT and reduced-model Monte Carlo diagnostics. Panel~(a) should therefore be interpreted as a campaign-scale local sensitivity result, whereas panel~(b) is a reduced-model robustness diagnostic for finite-horizon launch phase-out rather than a full-model Monte Carlo validation of architecture rankings.

\begin{figure}[htbp]
	\centering
	\begin{subfigure}[t]{0.485\linewidth}
		\centering
		\includegraphics[width=\linewidth]{figure/Fig6_Tornado_Sensitivity.png}
		\caption{OAT tornado sensitivity of reported campaign-scale metrics under the prescribed local perturbations.}
		\label{fig:tornado_sub}
	\end{subfigure}\hfill
	\begin{subfigure}[t]{0.485\linewidth}
		\centering
		\includegraphics[width=\linewidth]{figure/fig5_MonteCarlo_Robustness_A.png}
		\caption{Reduced-model Monte Carlo robustness diagnostic under joint perturbations. In the archived run, no finite-horizon phase-out was detected among numerically accepted samples, so the panel reports absence of detection rather than a non-empty year distribution.}
		\label{fig:mc_sub}
	\end{subfigure}
	\caption{Sensitivity and robustness diagnostics.
		Panel~(a) identifies the dominant local drivers of $T_b$, CI, and the Earth-launch burden proxy $J_E$ in the campaign-scale model.
		Panel~(b) summarizes the reduced-model Monte Carlo robustness diagnostic under joint perturbations around the adopted mature-regime baseline.}
	\label{fig:sens_tornado_mc}
\end{figure}

The OAT analysis identifies payload-per-launch \(L_R\) as the dominant local driver of build time, cost, and Earth-launch burden near the adopted mixed baseline comparator. The Earth-launch unit-cost coefficient \(C_E\) strongly affects the cost ranking while leaving build time essentially unchanged, whereas local perturbations in \(K_{\max,\mathrm{total}}\), \(\delta_{\mathrm{eff}}\), \(\alpha\), and \(r\) induce comparatively small variations around the same benchmark. The campaign-scale bridge sensitivity study supports the same interpretation on a wider full-model grid: across representative baseline and disruption cases, \(L_R\) remains the strongest driver, \(C_E\) primarily affects cost, and the explored variations in \(K_{\max,\mathrm{total}}\) and \(\delta_{\mathrm{eff}}\) do not move the system out of the rocket-dominant regime. Taken together, these campaign-scale diagnostics indicate that the baseline fast-build neighborhood is not only locally rocket-dominant, but also remains rocket-dominant across a broader set of key throughput-, payload-, and cost-side perturbations.

The reduced Monte Carlo diagnostic supports the same interpretation from a different angle. In the archived run, 25 samples satisfied the relaxed terminal-mass screen, but only 13 were numerically accepted under the reduced-model acceptance rule, and none of these accepted samples exhibited finite-horizon Earth\(\to\)Moon launch phase-out. Panel~(b) should therefore be read as an absence-of-detection result rather than as an empirical distribution of phase-out years. Within the adopted reduced baseline neighborhood and 50-year horizon, the sampled reduced regime remained in a no-phaseout / rocket-persistent region rather than in the vicinity of a robust finite-horizon phase-out boundary.

\subsection{Reduced mature-regime domain scan}
\label{subsec:phaseout_heatmap}

To test whether the no-phaseout behavior observed near the reduced baseline extends beyond the immediate Monte Carlo neighborhood, we next examine a reduced mature-regime two-dimensional diagnostic scan. In this archived layer, the reduced model is evaluated with \(\alpha=1\), finite disruption windows deactivated, and a fixed horizon of \(T=50\) years. Effective throughput is varied by scaling \(K_{\max,\mathrm{total}}\) while keeping \(\delta_{\mathrm{eff}}\) fixed, and Earth-launch-burden targets are generated as multipliers of a feasible reduced-model reference burden \(J_E^{\mathrm{ref}}\). The scan is therefore diagnostic rather than predictive: it is designed to test whether finite-horizon launch phase-out emerges anywhere in the explored reduced domain.

Table~\ref{tab:phaseout_scan_summary} summarizes the archived reduced-domain scan. The explored grid comprises five effective-throughput factors and five burden-target multipliers, for a total of 25 grid points. Across this scanned domain, all points were labeled ``Never'' for finite-horizon phase-out within the 50-year horizon. Most grid points were classified as approximate no-sign-change cases, indicating that the refinement procedure did not identify a finite-horizon transition even at the diagnostic level.

\begin{table}[htbp]
	\centering
	\caption{Summary of the archived reduced mature-regime domain scan.}
	\label{tab:phaseout_scan_summary}
	\small
	\begin{tabular}{ll}
		\hline
		Quantity & Archived result \\
		\hline
		Model layer & Reduced mature-regime diagnostic \\
		Horizon & \(T=50\) yr \\
		Commissioning setting & \(\alpha=1\) \\
		Finite disruption windows & Deactivated \\
		Effective-throughput factors & \(\{0.50,\,0.75,\,1.00,\,1.25,\,1.50\}\) \\
		Burden-target multipliers & \(\{0.1,\,0.2,\,0.5,\,1.0,\,2.0\}\) \\
		Total grid points & 25 \\
		Grid points with detected finite-horizon phase-out & 0 \\
		Grid points labeled ``Never'' & 25 \\
		Status counts & 24 \texttt{approx\_grid\_no\_sign\_change}, 1 \texttt{ok\_grid} \\
		\hline
	\end{tabular}
\end{table}

The reduced-domain scan therefore does not reveal a finite-horizon phase-out boundary within the explored range. Instead, it indicates no-phaseout persistence across the scanned reduced domain. This result does not prove that finite-horizon phase-out is impossible globally; rather, it shows that the presently explored mature-regime domain is insufficient to produce such a transition. In combination with the Monte Carlo diagnostic of Section~\ref{subsec:tornado_mc}, the scan supports a consistent reduced-model message: under the adopted baseline neighborhood and currently explored reduced domain, finite-horizon launch phase-out is not robustly realized.

These reduced diagnostics do not replace the campaign-scale results reported in Section~\ref{sec:results}. Their role is narrower and more mechanistic: they test whether a finite-horizon transition emerges under simplified mature-regime conditions and clarify that, within the present reduced-domain evidence base, no-phaseout persistence remains the dominant outcome.

All sensitivity outputs are archived in machine-readable form so that the reported plots, robustness summaries, and reduced-domain scan statistics can be regenerated directly from the corresponding run ledgers.


\section{Discussion}
\label{sec:discussion}

The results support a consistent interpretation of lunar logistics architecture selection as a regime-dependent systems problem rather than as a simple ranking among transport options. Across the campaign-scale baseline benchmark, hybrid trade-off endpoints, disruption tests, and sensitivity analyses, the same qualitative pattern emerges: under strong schedule pressure and the adopted baseline capacity assumptions, the optimized solution remains rocket-dominant, whereas SE-assisted transport becomes more consequential only when the optimization is allowed to trade schedule performance for lower Earth-launch burden.

\subsection{Why the baseline fast-build regime remains rocket-dominant}
\label{subsec:discussion_baseline}

The baseline fast-build regime is not rocket-dominant because the SE pathway is intrinsically infeasible, but because it is not sufficiently competitive under the combined requirements of rapid completion, finite transshipment buffering, and reliability-scaled effective throughput. Under the adopted launch-cadence cap, the direct rocket leg can sustain a build schedule close to the lower bound implied by aggregate Earth$\to$Moon launch capacity, whereas the SE-assisted chain remains constrained by commissioning, buffer-mediated staging, and reduced effective availability. When rapid completion is prioritized, the optimizer therefore has little incentive to route substantial mass through the elevator leg.

This same mechanism explains why the fast hybrid solution collapses toward the rocket-only benchmark. In principle, the hybrid architecture enlarges the admissible policy space and could exploit both pathways simultaneously. Under the adopted fast-build benchmark, however, the additional pathway does not create a materially competitive operating regime. Instead, the hybrid solution converges toward the rocket-dominant limit because the marginal contribution of the SE leg is too small to change the active bottleneck. The presence of an additional transport mode is therefore not, by itself, sufficient to guarantee meaningful utilization in an optimized campaign-scale schedule.

This interpretation is reinforced by the archived campaign-scale bridge sensitivity study. Across the representative cases \texttt{mixed\_base}, \texttt{mixed\_sefail}, \texttt{mixed\_rocketfail}, and \texttt{mixed\_bothfail}, the full model was re-solved under \(\tau_\star\)-oriented and cost-oriented policies while varying \(K_{\max,\mathrm{total}}\), \(\delta_{\mathrm{eff}}\), \(L_R\), \(\gamma_{\mathrm{env}}\), and \(C_E\) over multiplicative factors \(\{0.6,0.8,1.0,1.2,1.5\}\). Across all archived bridge runs, the regime label remained rocket-dominant. The strongest variations were induced by \(L_R\), whereas the explored changes in \(K_{\max,\mathrm{total}}\) and \(\delta_{\mathrm{eff}}\) produced comparatively small effects and did not trigger a regime flip. The baseline fast-build neighborhood is therefore not merely locally rocket-dominant at a single benchmark point, but remains rocket-dominant across a wider set of campaign-scale throughput-, payload-, and cost-side perturbations.

From a planning perspective, the relevant question is not whether an SE pathway can be included in the architecture, but whether it becomes active enough to change the schedule-defining bottleneck. Under the baseline assumptions examined here, that condition is not satisfied in the fast-build regime.

\subsection{Reliability, trade-off relaxation, and no-phaseout persistence}
\label{subsec:discussion_regime}

A central conclusion of the study is that the transition away from rocket-dominant transport is not reducible to a single scalar cutoff. Instead, architecture competitiveness depends jointly on effective SE throughput,
\[
K_{\mathrm{SE}}^{\mathrm{eff}}=\bar{\eta}K_{\mathrm{SE}}^{\mathrm{sys}},
\]
launch-side constraints, finite buffering, and schedule tolerance. The comparison between architectures is therefore better understood as a regime-identification problem in a coupled parameter space than as a one-threshold problem in throughput or launch aversion alone.

The campaign-scale disruption results reinforce this interpretation. When the optimized solution already places negligible mass on the SE leg, reductions in SE performance have little effect on the reported schedule or launch burden because they do not alter the active bottleneck. By contrast, rocket-side disruption changes the reported schedule precisely because the rocket leg remains decisive in the baseline fast-build regime. Under longer horizons or stronger Earth-launch-burden emphasis, the same logic implies that the SE leg becomes more consequential and the optimized solution can move toward mixed or SE-assisted operation along the feasible trade-off set.

The reduced-model diagnostics support the same message from a different angle, but with an important qualification. The reduced Monte Carlo experiment does not show a non-empty empirical distribution of finite-horizon phase-out years; instead, it indicates that no finite-horizon Earth$\to$Moon launch phase-out was detected among numerically accepted samples in the adopted mature-regime baseline neighborhood. Likewise, the archived reduced two-dimensional scan over effective-throughput and Earth-launch-burden conditions does not reveal a finite-horizon phase-out boundary within the explored domain. The most useful interpretation is therefore not that a space elevator becomes ``good'' beyond a single threshold, but that the presently explored baseline neighborhood and reduced domain remain in a no-phaseout / rocket-persistent region. Within the current evidence base, reduced diagnostics support no-phaseout persistence rather than robust finite-horizon transition discovery.

\subsection{Numerical credibility, modeling scope, and limitations}
\label{subsec:discussion_scope}

The numerical conclusions are supported by internal consistency audits and repeated-solve checks. Representative runs satisfy the construction and post-construction feasibility conditions with very small scaled constraint violations, the cumulative delivered-mass state remains monotone non-decreasing, and reruns under perturbed initial guesses preserve the same regime-level conclusions. In the archived representative campaign-scale consistency audit, all eight reruns were feasible, and the relative spreads in build time and total cost were \(3.18\times 10^{-13}\) and \(1.53\times 10^{-6}\), respectively, both far below the predefined 3\% tolerance. These results indicate that the reported campaign-scale benchmark is numerically stable under the adopted initialization perturbations.

These checks, however, establish numerical consistency rather than engineering validation. The present framework does not validate the real-world feasibility of a lunar space elevator, the completeness of the adopted Earth-launch-burden proxy, or the economic realism of the relative cost coefficients. Those quantities remain scenario assumptions used to support comparative architecture analysis at the planning scale.

The model is also intentionally aggregate. Launches are relaxed to planning-scale rates, Apex$\to$Moon launches are excluded from the Earth-launch-burden metric, and port-level heterogeneity is pooled into scenario parameters. These simplifications are appropriate for auditable architecture comparison and regime interpretation, but they also define the scope of the conclusions. In addition, the reduced-model diagnostics reported here do not establish a global absence of finite-horizon phase-out; they show only that such a transition was not detected in the adopted baseline neighborhood or within the currently explored reduced domain. Future work could therefore extend the reduced-domain scan, enrich the environmental indicators, incorporate explicit heterogeneity across ports or bases, and examine more detailed disruption-and-recovery dynamics while preserving the same underlying mass-balance structure.

% =========================================================
\section{Conclusion}

This study developed a unit-consistent, reliability-aware mass-flow framework for comparing rocket-only, SE-assisted two-stage, and hybrid Earth--Moon logistics architectures under shared launch-rate caps, finite buffering, throughput ramp-up, and post-construction resupply requirements. The analysis combined campaign-scale architecture comparison, numerical trade-off construction, disruption testing, local sensitivity analysis, campaign-scale bridge sensitivity, repeated-solve consistency audits, and reduced regime-oriented diagnostics to determine when SE-assisted transport remains inactive, becomes more relevant, or becomes structurally important.

Under the baseline benchmark, rocket-only transport achieves the shortest build schedule (\(T_b\approx 8.01\) years) but does so at the highest Earth-launch burden in the adopted metric. By contrast, the SE-assisted two-stage architecture sharply reduces the reported Earth-launch burden, but only by accepting a much longer construction horizon (\(T_b\approx 235.91\) years). At the same fast-build benchmark, the hybrid solution remains effectively rocket-dominant, with limited SE utilization and correspondingly weak sensitivity to SE-side disruption. The mere inclusion of an SE leg in the architecture therefore does not guarantee meaningful use in the optimized schedule.

More broadly, the results support a regime-based interpretation of lunar logistics planning. The practically relevant output is not a single preferred architecture, but a characterization of the operating regions in which rocket-only, hybrid, or SE-assisted transport becomes competitive. Under the adopted parameterization, the baseline fast-build region is rocket-dominant, and movement toward more substantial SE participation occurs only when schedule pressure is relaxed and Earth-launch burden is weighted more strongly in the optimization. The archived campaign-scale bridge sensitivity study strengthens this interpretation by showing that, across representative baseline and disruption cases and across a wider grid of throughput-, payload-, and cost-side perturbations, the explored full-model solutions remain rocket-dominant rather than crossing into a distinct SE-led regime.

The archived reduced-model diagnostics support the same message from a narrower mechanism-oriented perspective. In the adopted mature-regime baseline neighborhood, no finite-horizon Earth--Moon launch phase-out was detected among numerically accepted Monte Carlo samples, and the explored reduced two-dimensional domain likewise did not reveal a finite-horizon phase-out boundary. Within the present evidence base, the reduced diagnostics therefore support no-phaseout persistence rather than a robust finite-horizon transition.

Within the scope of the adopted stress-test setting and the present planning-scale aggregate model, the numerical evidence reported here, together with the archived trajectories, sensitivity batches, and audit logs described in the availability statements, indicates that these conclusions are stable at the regime level under local perturbations, wider campaign-scale bridge perturbations, and reruns with perturbed initial guesses. Future work can refine the framework with richer environmental indicators, explicit heterogeneity across ports or bases, more detailed failure-and-recovery dynamics, and a wider reduced-domain scan. Even with these limitations, the present results already show that reliability degradation, conservative effective throughput, finite buffering, and strong schedule pressure can be sufficient to keep SE-dominant logistics outside the practically competitive regime.


\section*{Data availability}
The computational artifacts supporting the conclusions of this study include campaign-scale scenario summaries, representative time-series trajectories, consistency-audit ledgers, campaign-scale sensitivity outputs, reduced-model Monte Carlo robustness logs, and reduced mature-regime domain-scan records. These materials are archived with the study data package and are sufficient to regenerate the reported tables and figures at the level of the published analysis. The public repository DOI or persistent URL must be inserted here after final archival deposition and before submission.

\section*{Code availability}
The archived code package includes the campaign-scale scenario runner, the single-case runner, the consistency-audit workflow, the campaign-scale OAT sensitivity workflow, the reduced-model Monte Carlo robustness workflow, the reduced mature-regime domain-scan script, and the automated paper-generation pipeline used to regenerate the principal computational artifacts. The public repository DOI or persistent URL must be inserted here after final archival deposition and before submission, together with environment information and minimal reproduction commands.

\section*{Author contributions}
Y.W. contributed to conceptualization, methodology, software, formal analysis, and writing of the original draft. Q.L. contributed to methodology, validation, scenario formulation, and manuscript review and editing. J.H. contributed to data curation, post-processing, sensitivity analysis, robustness analysis, and manuscript review and editing. All authors reviewed and approved the final manuscript.

\section*{Competing interests}
The authors declare no competing interests.

\begin{thebibliography}{99}
	
	\bibitem{COMAP2026B}
	COMAP.
	2026 MCM Problem B: Creating a Moon Colony Using a Space Elevator System.
	\textit{Problem statement (PDF)}
	\url{https://www.contest.comap.com/undergraduate/contests/mcm/contests/2026/problems/2026_MCM_Problem_B.pdf} (2026).
	
	\bibitem{Smitherman2000SpaceElevators}
	Smitherman, D. V.
	\textit{Space Elevators: An Advanced Earth-Space Infrastructure for the New Millennium}.
	NASA/CP-2000-210429 (2000).
	
	\bibitem{Edwards2000NIAC}
	Edwards, B. C.
	The Space Elevator NIAC Phase II Final Report.
	\textit{NIAC Final Report}
	\url{https://www.niac.usra.edu/files/studies/final_report/521Edwards.pdf} (2003).
	
	\bibitem{Chang2022}
	Barry, K. M. \& Alfaro, E. P.
	Changing the economic paradigm for building a space elevator.
	\textit{Acta Astronaut.} \textbf{200}, 586--598 (2022).
	doi:10.1016/j.actaastro.2022.08.025.
	
	\bibitem{Nixon2023TetherMaterials}
	Nixon, A., Knapman, J. \& Wright, D. H.
	Space elevator tether materials: An overview of the current candidates.
	\textit{Acta Astronaut.} \textbf{210}, 483--487 (2023).
	doi:10.1016/j.actaastro.2023.04.008.
	
	\bibitem{Swan2023ApexAnchorInitialResearch}
	Swan, P. A.
	Apex Anchor initial research.
	\textit{IAC/IAF paper (ISEC archive)}
	\url{https://space-elevator.squarespace.com/s/ISEC-2024-IAC-Apex-Anchor-Initial-Research-paper.pdf} (2023).
	
	\bibitem{Swan2025GalacticHarbours}
	Swan, P. \& Griffin, D.
	Galactic harbours define permanent space access transportation infrastructures.
	\textit{International Astronautical Congress (IAC)}, Sydney, Australia.
	\url{https://www.isec.org/s/ISEC-2025-IAC-Galactic-Harbours-Define-Permanent-Space-Access-Infrastructures-paper.pdf} (2025).
	
	\bibitem{Merancy2025ADD}
	Merancy, N. \textit{et al.}
	\textit{ESDMD-001: Moon to Mars Architecture Definition Document, Revision C}.
	NASA/TM-20250010956.
	\url{https://ntrs.nasa.gov/api/citations/20250010956/downloads/ADD%20Revision%20C%20-%2020260115.pdf?attachment=true} (2025).
	
	\bibitem{Azami2024LBMC}
	Azami, M., Kazemi, Z., Moazen, S., Dub\'e, M., Potvin, M.-J. \& Skonieczny, K.
	A comprehensive review of lunar-based manufacturing and construction.
	\textit{Prog. Aerosp. Sci.} \textbf{150}, 101045 (2024).
	doi:10.1016/j.paerosci.2024.101045.
	
	\bibitem{Xia2025LunarBaseInfra}
	Xia, Y., Gao, Y., Han, W., Li, X., Zhou, C., Zhou, Y. \& Ding, L.
	Lunar base infrastructure construction: Challenges and future directions.
	\textit{Autom. Constr.} \textbf{176}, 106251 (2025).
	doi:10.1016/j.autcon.2025.106251.
	
	\bibitem{Vasudeva2026LunarConstruction}
	Vasudeva, K., Emami, M. R. \& Dickinson, C.
	Lunar construction: A state-of-the-art survey.
	\textit{Prog. Aerosp. Sci.} \textbf{161}, 101159 (2026).
	doi:10.1016/j.paerosci.2025.101159.
	
	\bibitem{Li2024LaunchLanding}
	Li, X., Gao, Y., Zhou, Y., Han, W. \& Zhou, C.
	A review on design and construction of the lunar launch/landing infrastructure.
	\textit{Adv. Space Res.} \textbf{74}, 4030--4049 (2024).
	doi:10.1016/j.asr.2024.07.006.
	
	\bibitem{Ishimatsu2016EMMLogistics}
	Ishimatsu, T., de Weck, O. L., Hoffman, J. A., Ohkami, Y. \& Shishko, R.
	Generalized multicommodity network flow model for the Earth--Moon--Mars logistics system.
	\textit{J. Spacecr. Rockets} \textbf{53}(1), 25--38 (2016).
	doi:10.2514/1.A33235.
	
	\bibitem{Ho2016FlexiblePath}
	Ho, K., de Weck, O. L., Hoffman, J. A. \& Shishko, R.
	Campaign-level dynamic network modelling for spaceflight logistics for the flexible path concept.
	\textit{Acta Astronaut.} \textbf{123}, 51--61 (2016).
	doi:10.1016/j.actaastro.2016.03.006.
	
	\bibitem{Ho2024Review}
	Ho, K.
	Space logistics modeling and optimization: Review of the state of the art.
	\textit{J. Spacecr. Rockets} \textbf{61}(5), 1417--1427 (2024).
	doi:10.2514/1.A35982.
	
	\bibitem{Apa2025EarthMoonLogistics}
	Apa, R., Hudson, J. \& Romano, M.
	Request-centric modeling and architecture optimization for Earth--Moon space logistics.
	\textit{J. Spacecr. Rockets} (2025).
	doi:10.2514/1.A36495.
	
	\bibitem{Ryan2022}
	Ryan, R. G., Marais, E. A., Balhatchet, C. J. \& Eastham, S. D.
	Impact of rocket launch and space debris air pollutant emissions on stratospheric ozone and global climate.
	\textit{Earth's Future} \textbf{10}, e2021EF002612 (2022).
	doi:10.1029/2021EF002612.
	
	\bibitem{Brown2023SustainableLaunch}
	Brown, T. F. M., Bannister, M. T. \& Revell, L. E.
	Envisioning a sustainable future for space launches: A review of current research and policy.
	\textit{J. R. Soc. N. Z.} \textbf{54}(3), 273--289 (2024).
	doi:10.1080/03036758.2022.2152467.
	
	\bibitem{Sirieys2022AtmosphericImpact}
	Sirieys, E., Gentgen, C., Jain, A., Milton, J. \& de Weck, O. L.
	Space sustainability isn’t just about space debris: On the atmospheric impact of space launches.
	\textit{MIT Science Policy Review} (2022).
	\url{https://sciencepolicyreview.org/wp-content/uploads/securepdfs/2022/08/MITSPR-v3-191618003013.pdf}.
	
	\bibitem{Barker2024Inventory}
	Barker, C. R., Marais, E. A. \& McDowell, J.
	Global 3D rocket launch and re-entry air pollutant and CO$_2$ emissions at the onset of the megaconstellation era.
	\textit{Sci. Data} \textbf{11}, 1079 (2024).
	doi:10.1038/s41597-024-03910-z.
	
	\bibitem{Sharma2024Atmosphere}
	Sharma, S. P.
	\textit{Impact of Spaceflight on Earth's Atmosphere}.
	NASA/TM-20240013276.
	\url{https://ntrs.nasa.gov/api/citations/20240013276/downloads/NASA-TM-20240013276-V6.pdf} (2024).
	
	\bibitem{Revell2025Ozone}
	Revell, L. E. \textit{et al.}
	Near-future rocket launches could slow ozone recovery.
	\textit{npj Clim. Atmos. Sci.} \textbf{8}, 212 (2025).
	doi:10.1038/s41612-025-01098-6.
	
	\bibitem{Ewert2022BVAD}
	Ewert, M. K. \textit{et al.}
	Life Support Baseline Values and Assumptions Document (BVAD).
	\textit{NASA Technical Reports Server}
	\url{https://ntrs.nasa.gov/api/citations/20210024855/downloads/BVAD_2.15.22-final.pdf} (2022).
	
	\bibitem{NASASTD3001V2RevE2025}
	NASA.
	NASA-STD-3001, Volume 2, Revision E: Spaceflight Human-System Standard.
	\url{https://standards.nasa.gov/system/files/tmp/NASA-STD-3001\%20Vol\%202\%20Rev\%20E_1.pdf} (2025).
	
	\bibitem{NASA2023WaterHumanConsumptionTB}
	NASA Office of the Chief Health \& Medical Officer (OCHMO).
	OCHMO-TB-027: Water -- Human Consumption (NASA-STD-3001 Technical Brief).
	\url{https://www.nasa.gov/wp-content/uploads/2023/12/ochmo-tb-027-water.pdf} (2023).
	
	\bibitem{NASA2023ECLSSBrief}
	NASA Office of the Chief Health \& Medical Officer (OCHMO).
	OCHMO-TB-002: Environmental Control \& Life Support System (ECLSS): Human-Centered Approach (NASA-STD-3001 Technical Brief).
	\url{https://www.nasa.gov/wp-content/uploads/2023/07/eclss-technical-brief-ochmo.pdf} (2023).
	
	\bibitem{Williamson2023ISSWater}
	Williamson, J. \textit{et al.}
	Status of ISS Water Management and Recovery.
	\textit{International Conference on Environmental Systems (ICES)}
	\url{https://ntrs.nasa.gov/api/citations/20230006217/downloads/ICES\%202023-097\%20Status\%20of\%20ISS\%20Water\%20Management\%20and%20Recovery.pdf} (2023).
	
	\bibitem{Norris1998MarkovChains}
	Norris, J. R.
	\textit{Markov Chains}
	(Cambridge University Press, 1998).
	
	\bibitem{Kirk2004OptimalControl}
	Kirk, D. E.
	\textit{Optimal Control Theory: An Introduction}
	(Dover Publications, 2004).
	
	\bibitem{BrysonHo1975OptimalControl}
	Bryson, A. E. \& Ho, Y.-C.
	\textit{Applied Optimal Control}
	(Taylor \& Francis, 1975).
	
	\bibitem{AriasMontano2015}
	Arias-Monta\~no, A., Coello Coello, C. A. \& Mezura-Montes, E.
	Multiobjective evolutionary algorithms in aeronautical and aerospace engineering.
	\textit{IEEE Trans. Evol. Comput.} \textbf{16}(5), 662--694 (2012).
	doi:10.1109/TEVC.2011.2169968.
	
	\bibitem{SpaceX2025FalconUsersGuide}
	SpaceX.
	\textit{Falcon Payload User's Guide} (Version 2025-05-09).
	\url{https://www.spacex.com/assets/media/falcon-users-guide-2025-05-09.pdf} (2025).
	
	\bibitem{NASA2023SLSFactsheet}
	NASA.
	Space Launch System (SLS) Fact Sheet.
	\url{https://www.nasa.gov/wp-content/uploads/2023/07/0080-sls-fact-sheet-jun2023-508-1.pdf} (2023).
	
\end{thebibliography}

\appendix

\section{Supplementary Computational and Modeling Details}
\label{app:supplementary_methods}

This appendix provides supplementary methodological and computational details that support the main text but are not required to follow the principal scientific argument. The material is organized around five themes: computational conventions, the planning-scale interpretation of reliability, supplementary details of the numerical workflow for the campaign-scale and reduced-model calculations, archived outputs and audit logic, and additional notes for reported metrics and reduced-model phase-out diagnostics.

\subsection{Computational conventions and planning-scale implementability}
\label{app:comp_conventions}

All model quantities are defined in physical units: mass in metric tons (t), mass flow in t/year, and launch rates in launches/year. For numerical conditioning, plotting convenience, or ledger readability, the implementation may apply fixed presentation scalings. These scalings do not alter the meaning of the variables, constraints, or reported feasibility conditions. In particular, labels such as ``kLaunch/yr'' denote only a display scaling by $10^3$, and all archived ledgers, trajectory files, and exported summary tables remain stored in physical units with explicit unit headers.

The campaign-scale optimization is formulated using bounded launch-rate controls rather than integer launch counts. This relaxation is justified because the construction target implies
\[
\frac{M_{\mathrm{goal}}}{L_R}\sim 10^6,
\]
so the planning problem involves sufficiently large launch counts for rate-based control to provide an accurate campaign-scale approximation.

To bound the approximation error, suppose that over a discretization interval of length $\Delta t$, the optimized launch-rate controls $F_{EM,k}$ and $F_{AM,k}$ are converted to implementable integer launch counts by nearest-integer rounding of $\Delta t F_{EM,k}$ and $\Delta t F_{AM,k}$. Then the count error per interval is at most $1/2$ launch per leg. Hence the absolute cumulative delivered-mass error over $K$ periods is bounded by
\[
\frac{1}{2} K L_R
\]
per rocket leg, or by the corresponding two-leg bound when both pathways are rounded independently. Relative to the construction target, the delivery error therefore satisfies
\[
\frac{\text{cumulative mass error}}{M_{\mathrm{goal}}}
\le
\frac{\tfrac12 K L_R}{M_{\mathrm{goal}}},
\]
which is negligible at the campaign scale considered here. This supports the use of continuous launch-rate controls for architecture comparison, frontier construction, and regime identification.

\subsection{Planning-scale interpretation of the mean availability factor}
\label{app:ctmc_eta}

In the main text, imperfect SE operation is represented through a mean availability factor $\bar{\eta}\in[0,1]$ that scales effective throughput. This quantity may be interpreted as the stationary mean availability of a finite-state continuous-time Markov chain (CTMC), for example with operational, degraded, and failed modes \cite{Norris1998MarkovChains}.

Let the SE operating mode be a finite-state CTMC with state space
\[
\mathcal{S}=\{1,\dots,m\},
\]
generator matrix $Q$, stationary distribution $\boldsymbol{\pi}$, and state-dependent throughput multipliers
\[
\gamma_i\in[0,1],\qquad i\in\mathcal{S}.
\]
If $K_{\mathrm{SE}}^{\mathrm{sys}}$ denotes the nominal system throughput under perfect operation, then the long-run effective throughput is
\[
K_{\mathrm{SE}}^{\mathrm{eff}}
=
\left(\sum_{i=1}^m \pi_i \gamma_i\right)K_{\mathrm{SE}}^{\mathrm{sys}}.
\]
Accordingly,
\[
\bar{\eta}:=\sum_{i=1}^m \pi_i \gamma_i .
\]

The optimization model depends only on the aggregate factor $\bar{\eta}$, not on the detailed transition structure of the CTMC. This reduction is appropriate when the planning horizon is long relative to the mixing time of the mode-switching process and when availability statistics are stable at the annual scale. Under these conditions, long-run average throughput is the dominant quantity for campaign-scale feasibility and trade-off analysis.

For resilience stress tests, the effective-capacity bound may additionally be multiplied by a deterministic time-window health factor. This extension represents finite disruption windows without changing the structure of the optimization problem. The resulting experiments should therefore be interpreted as timing-based stress tests rather than as full stochastic pathwise simulations of failure and repair.

\subsection{Supplementary note on post-construction water resupply}
\label{app:water_extra}

To connect construction-phase logistics with near-term operational sustainment, the model imposes a one-year post-construction water-resupply requirement for a colony of
\[
P=100{,}000
\]
residents \cite{COMAP2026B}. Gross daily water use is bounded by
\[
w_{\mathrm{gross}}\in[3.0,4.5]\ \text{kg/(person$\cdot$day)},
\]
and the closed-loop recovery fraction is taken in the range
\[
\rho\in[0.90,0.98].
\]
The corresponding net annual resupply mass is
\[
W_{\mathrm{net}}=(1-\rho)\,P\,w_{\mathrm{gross}}\frac{365}{1000}
\quad \text{t/year}.
\]

Under the baseline values $(w_{\mathrm{gross}},\rho)=(3.75,\,0.95)$, this gives
\[
W_{\mathrm{net}}=6843.8\ \text{t/year}.
\]
Although this resupply requirement is small relative to the total construction target, it is retained to ensure that a construction-feasible policy also remains admissible under immediate post-build operational demand.

Additional consumables beyond water are not included in the baseline formulation. Instead, they are considered only in sensitivity analyses using traceable bounds from life-support references \cite{Ewert2022BVAD}. This separation prevents secondary sustainment assumptions from dominating the primary comparison of transport architectures while still allowing robustness checks under expanded operational loads.

\subsection{Supplementary numerical workflow for campaign-scale runs and reduced-model diagnostics}
\label{app:numerics}

Two computational layers are used. First, campaign-scale baseline comparisons are generated under the unified construction-plus-one-year-resupply model introduced in the main text. Second, a reduced mature-regime formulation is used for mechanism-oriented diagnostics of no-phaseout persistence and possible finite-horizon launch transitions. Both layers share the same physical interpretation of flows, buffer coupling, and Earth-launch-burden metrics, but they serve different purposes: the former quantifies architecture-level schedule and cost trade-offs, whereas the latter tests whether threshold-like transition behavior emerges under mature operating conditions.

Time is discretized using a uniform step size $\Delta t$ (years), with monthly planning corresponding to $\Delta t=1/12$. Piecewise-constant controls are used on each interval $[t_k,t_{k+1})$, where $t_k=k\Delta t$. In the campaign-scale model, the build phase and the subsequent one-year resupply phase are represented on separate grids, but both use the same physical units and flow definitions.

The discretized dynamics follow directly from the continuous flow model:
\begin{align}
	A_{k+1} &= A_k + \Delta t\left(\Phi_{SE,k}-\Phi_{AM,k}\right), \label{eq:disc_A_rev_app}\\
	M_{k+1} &= M_k + \Delta t\left(\Phi_{EM,k}+\Phi_{AM,k}\right). \label{eq:disc_M_rev_app}
\end{align}
The Apex buffer constraint is enforced at every step through
\[
0\le A_k\le A_{\max}.
\]
Hence feasibility depends not only on average throughput but also on transient synchronization between elevator inflow and downstream Apex-to-Moon launches.

For a prescribed construction deadline $\bar{T}$, define
\[
K_b=\left\lceil \bar{T}/\Delta t \right\rceil,\qquad
K_w=\left\lceil 1/\Delta t \right\rceil .
\]
Completion by the deadline is enforced by
\begin{equation}
	M_{K_b}\ge M_{\mathrm{goal}},
	\label{eq:disc_completion_by_deadline_app}
\end{equation}
and the one-year post-construction resupply requirement is enforced by
\begin{equation}
	M_{K_b+K_w}-M_{K_b}\ge W_{\mathrm{net}}.
	\label{eq:disc_water_by_deadline_app}
\end{equation}
For reporting, the achieved completion time on the grid is defined by
\[
T_{\mathrm{rep}}:=\min\{k\Delta t:\ M_k\ge M_{\mathrm{goal}}\}\le K_b\Delta t.
\]
This formulation is conservative because it guarantees a full one-year delivery window after the imposed deadline even if the mass target is reached earlier within the optimized trajectory.

In the campaign-scale model, the effective elevator capacity is represented through the sampled bound
\begin{equation}
	\Phi_{SE,k}^{\max}=N_{SE,k}\,\eta_k\,K_{\mathrm{SE}}^{\mathrm{sys}},
	\label{eq:phi_se_max_k_app}
\end{equation}
where $N_{SE,k}=N_{SE}(t_k)$ is the discretized commissioning factor and $\eta_k$ is the availability factor. For the main architecture comparisons, the mean-availability formulation is used,
\[
\eta_k\equiv\bar{\eta}\qquad \text{for all }k.
\]
For resilience illustrations and selected sensitivity tests, structured time windows are imposed in which elevator or rocket capacity is reduced. In the reduced diagnostic model, early commissioning transients and finite disruption windows are removed so that the influence of effective throughput and Earth-launch-burden conditions can be examined more directly.

For each baseline scenario and each strategy class, three representative policy variants are computed: a minimum-feasible-build-time solution $\tau_\star$ obtained by bisection on the build horizon, a schedule-oriented solution near $\tau_\star$, and a cost-oriented solution constrained to remain within a small relative tolerance of $\tau_\star$. In the baseline runs, this tolerance is set to
\[
\tau \le (1+\epsilon_\tau)\tau_\star,\qquad \epsilon_\tau=0.05.
\]
This construction preserves near-minimal schedule length while allowing comparison of lower-cost feasible operating patterns within the same scenario.

Across the campaign-scale runs, feasible solutions from the three strategy classes are compared using the reported objective triple
\[
(T_{\mathrm{rep}},\mathrm{CI},\mathrm{LI}),
\]
where $\mathrm{LI}=J_E/J_E^{\mathrm{ref}}$ is the normalized launch-intensity index. A solution $\mathbf{J}^{(1)}$ is treated as dominating $\mathbf{J}^{(2)}$ if $J_i^{(1)}\le J_i^{(2)}$ for all components and strict inequality holds in at least one component, subject to a small numerical tolerance. The filtered non-dominated set is then used to visualize architecture-dependent trade-offs in completion time, normalized cost, and normalized launch intensity.

The archived reduced diagnostics are of two kinds. First, a reduced Monte Carlo workflow perturbs \(\{K_{\max,\mathrm{total}},L_R,r,C_E,C_A\}\) independently by prescribed multiplicative factors around the adopted mature-regime baseline and records feasibility status, realized launch burden, and finite-horizon phase-out status. Second, a reduced two-dimensional domain scan varies effective-throughput factors and Earth-launch-burden target multipliers over a predefined grid. In the archived main-text results, these diagnostics are used to test whether finite-horizon launch phase-out emerges in the baseline neighborhood or explored reduced domain; they are not interpreted as substitutes for the campaign-scale architecture comparisons.

When a target launch-burden level is imposed in the reduced-domain scan, the corresponding control parameter is matched numerically rather than through a closed-form exact solve. A penalty or target parameter is varied over a logarithmic grid and, when a feasible bracket is available, may be refined numerically so that the realized launch-burden level approaches the target value. A sampled point is accepted when the relative mismatch between realized and target launch intensity falls within the prescribed tolerance; in the default runs, this relative tolerance is $5\%$.

In the reduced-domain scan, the launch-burden grid is scaled from a feasible reference run rather than from an unconstrained limiting case. This avoids assuming that a vanishing launch-burden penalty is feasible and yields a reproducible family of target levels tied to an actually solved trajectory.

\subsection{Archived outputs, audits, and reproducibility}
\label{app:reproducibility}

To support independent verification, the implementation exports machine-readable run records, representative time-series trajectories, and automated feasibility audits. For campaign-scale baseline runs, the archived summary record contains at least the scenario label, strategy class, solver status, completion time, cost metric, launch-intensity metrics, and the corresponding time-series file path. For reduced-model diagnostics, the archived record contains the sampled control coordinates together with realized objective quantities, relative mismatch, regime labels when applicable, and feasibility flags.

Representative trajectories include at least
\[
t,\ A(t),\ M(t),\ F_{EM}(t),\ F_{AM}(t),\ \Phi_{SE}(t).
\]
The automated audits check mass-balance residuals for the discretized state equations, Apex buffer bound compliance, nonnegativity of state and control variables, monotonicity of cumulative delivered mass $M(t)$, and consistency between integrated inflow histories and reported peak inventory levels.

The discretized problems are solved by direct transcription using CasADi with IPOPT. To improve numerical conditioning, the implementation uses scaled objective values and scaled constraint residuals in the nonlinear program while preserving physical units in all archived outputs. Numerical scaling therefore affects neither the feasibility logic nor the physical interpretation of the reported trajectories and metrics.

To assess numerical stability, randomized-initialization consistency checks are performed in which repeated solves are run with perturbed initial guesses and perturbed initial horizon guesses. The resulting summary records the relative spread in completion time and total cost across successful runs and flags whether the prescribed consistency tolerances are met.

The archived output package additionally records the baseline parameter set, selected solver settings, and random seeds for restart-based checks. These archived artifacts are intended as computational verification aids rather than as additional optimization constraints. They do not alter the feasible set or the reported objective values, but they support numerical transparency and consistency checking of the reported results.

\section{Supplementary Notes for Reported Metrics and Resilience Diagnostics}
\label{app:results_notes}

\subsection{Reported cumulative quantities used in the baseline tables}
\label{app:reported_metrics}

For any feasible policy with reported completion time $T_{\mathrm{rep}}$, define the total reporting horizon as
\[
T_{\mathrm{tot}}:=T_{\mathrm{rep}}+1 .
\]
The cumulative quantities reported in Tables~\ref{tab:baseline_strategy_tau}--\ref{tab:mixed_tradeoff} are
\[
N_{EM}:=\int_{0}^{T_{\mathrm{tot}}}F_{EM}(t)\,dt,\qquad
N_{AM}:=\int_{0}^{T_{\mathrm{tot}}}F_{AM}(t)\,dt,
\]
\[
M_{SE}:=\int_{0}^{T_{\mathrm{tot}}}\Phi_{SE}(t)\,dt,\qquad
A_{\mathrm{peak}}:=\max_{t\in[0,T_{\mathrm{tot}}]}A(t).
\]
In the discrete implementation, these quantities are computed by the corresponding Riemann sums or sampled maxima.

\subsection{Auxiliary resilience metrics derived from archived trajectories}
\label{app:resilience_metrics}

Let the outage window be $[t_1,t_2]$, and let $M_{\mathrm{baseline}}(t)$ and $M_{\mathrm{shock}}(t)$ denote the lunar-delivered mass trajectories under the baseline and disrupted scenarios, respectively. The following auxiliary resilience metrics can be computed directly from the archived time-series outputs:
\[
\Delta M:=M_{\mathrm{baseline}}(t_2)-M_{\mathrm{shock}}(t_2),
\]
which measures the delivered-mass shortfall accumulated by the end of the outage window, together with the recovery time $t_{\mathrm{rec}}-t_2$ required to regain the baseline trajectory (or to reach $M_{\mathrm{goal}}$), and the compensation intensity measured by the peak and time-integrated excess of rocket cadence above baseline over a compensation window $[t_1,t_{\mathrm{rec}}]$.

These quantities are not used as primary optimization objectives in the main text, but they provide a compact diagnostic description of recovery severity across different outage timings and affected transport legs.

\subsection{Consistency of reported mass trajectories}
\label{app:mass_monotonicity}

Throughout the paper, $M(t)$ denotes cumulative lunar-delivered mass. Under the adopted dynamics and nonnegative controls, $M(t)$ is monotone non-decreasing. In continuous time, the mass-balance equation implies $\dot M(t)\ge 0$, and in discrete time the transcription implies $M_{k+1}\ge M_k$. A decreasing trajectory would therefore correspond to a different reported quantity, such as remaining unmet mass or a net-usable-mass state after additional deductions, and would require an explicit redefinition of the state variables. The figures and tables reported here remain consistent with the cumulative-delivered-mass interpretation throughout.

\subsection{\texorpdfstring{Phase-out rule for Earth $\to$ Moon launches}{Phase-out rule for Earth -> Moon launches}}
\label{app:phaseout_year}

This subsection defines the post-processing rule used to determine whether Earth$\to$Moon launches phase out within a reporting horizon. The rule is diagnostic only and is not imposed as an optimization constraint.

Let $\varepsilon_F>0$ be a small launch-rate threshold (launches/yr), and let $m\in\mathbb{N}$ be a confirmation-window length measured in time steps. To determine whether Earth$\to$Moon launches phase out within the reporting horizon, we first identify whether there exists an active-launch episode, defined as a contiguous block of at least $m$ time samples for which
\[
F_{EM}(t)>\varepsilon_F,
\]
after mild smoothing if needed to suppress discretization dithering.

If no such active block exists over the reporting horizon, phase-out is recorded as not detected within the horizon. Otherwise, let $t_s$ be the earliest time after the last detected active block such that
\[
\overline{F}_{EM}(t)\le \varepsilon_F^{\mathrm{rel}}
\qquad \text{for all } t\ge t_s,
\]
where $\overline{F}_{EM}(t)$ denotes a trailing moving average of $F_{EM}(t)$ and
\[
\varepsilon_F^{\mathrm{rel}}=\kappa_{\mathrm{rel}}\varepsilon_F,
\qquad \kappa_{\mathrm{rel}}\ge 1.
\]
The relaxation factor $\kappa_{\mathrm{rel}}$ is introduced to avoid false negatives caused by numerical dithering near zero. The values $(\varepsilon_F,m,\kappa_{\mathrm{rel}})$ used in each experiment are recorded in the run record.

This sustained-low criterion avoids declaring phase-out from isolated low-cadence intervals and is consistent with the audit-style detection logic used in the archived reduced-model post-processing analyses.

\subsection{Sampling conventions for OAT and Monte Carlo analyses}
\label{app:uq_sampling}

The OAT sensitivity study perturbs one parameter at a time around the campaign-scale baseline benchmark while holding the remaining parameters fixed. Unless a scenario-defined range is prescribed, the perturbation magnitude follows the same local convention as the archived sensitivity batch.

For the archived reduced Monte Carlo robustness diagnostic, selected reduced-model inputs are perturbed simultaneously by independent multiplicative factors around their baseline values. In the main-text diagnostic run, the perturbed set is \(\{K_{\max,\mathrm{total}},L_R,r,C_E,C_A\}\), the perturbation amplitude is \(\pm 5\%\), the requested launch-aversion parameter is \(\gamma=0.08\), and terminal delivered mass is screened using \(M_T\ge 0.99\,M_{\mathrm{goal}}\). For each sampled parameter vector, the reduced mature-regime problem is re-solved, and the resulting feasibility status, launch burden, and finite-horizon no-phaseout / phase-out status are archived together with the sampled parameters.

When randomized sampling is used, the run record stores the perturbation ranges, random seed, and sampled parameters for each run. Scenario-level outputs are archived in the same general machine-readable format as deterministic runs so that sensitivity plots, robustness summaries, and reduced-domain scan statistics can be regenerated directly from the stored records.

\section{Supplementary Numerical Credibility and Audit Notes}
\label{app:credibility}

\subsection{Automated consistency audit for representative runs}
\label{app:consistency_report}

To support internal validity and reproducibility, representative runs are audited for numerical feasibility, mass-balance consistency, monotonicity of the delivered-mass state, and agreement between archived trajectories and reported summary metrics. These checks are computed directly from archived run records and time-series outputs and are therefore independent of the narrative interpretation of any individual figure or table.

The archived campaign-scale audit verifies solver success, feasibility of the construction and one-year operational phases, late-window launch behavior, monotonicity of cumulative lunar-delivered mass, and consistency between archived time-series outputs and reported summary quantities. Reported constraint violations remain far below the practical feasibility thresholds used in the implementation, and the delivered-mass state remains monotone non-decreasing in the audited trajectories. Full machine-readable audit ledgers are included in the archived data package.

\subsection{Randomized-initialization consistency checks}
\label{app:audit_findings}

A common concern in optimization-based studies is that reported conclusions may depend strongly on solver initialization. Representative cases are therefore re-solved with perturbed initial guesses and perturbed initial horizon guesses, and the spread of the resulting completion times and total costs is summarized. This test is implemented through the archived repeated-solve workflow and provides a direct check on numerical stability under the same problem formulation.

\begin{table}[htbp]
	\centering
	\caption{Repeated-solve consistency summary for the representative campaign-scale mixed baseline benchmark.}
	\label{tab:consistency_summary_app}
	\small
	\begin{tabular}{lc}
		\hline
		Quantity & Archived value \\
		\hline
		Number of reruns & 8 \\
		All reruns feasible & Yes \\
		Feasible fraction & 1.0 \\
		Relative spread in build time & \(3.18\times 10^{-13}\) \\
		Relative spread in total cost & \(1.53\times 10^{-6}\) \\
		Build-time tolerance & 0.03 \\
		Cost tolerance & 0.03 \\
		Consistency pass & Yes \\
		\hline
	\end{tabular}
\end{table}

The relative spreads in build time and total cost are many orders of magnitude smaller than the prescribed tolerances. Small solver-path differences therefore do not alter feasibility classification or the reported regime-level conclusions for the representative campaign-scale benchmark.


\end{document}